\documentclass[12pt]{article}
\usepackage{a4wide}

\usepackage{amsmath,amssymb,amsthm}
\usepackage[mathscr]{eucal}
\usepackage{graphics}
\usepackage[hypertex]{hyperref}
\usepackage{epsfig}
\usepackage{psfrag}
\usepackage[usenames]{color}

%\renewcommand{\baselinestretch}{1.5}

%\definecolor{red}{rgb}{1,0.00,0.00}
%\definecolor{blue}{rgb}{0,0.08,0.55}
%\definecolor{green}{rgb}{0.0,0.6,0}
%\newcommand{\red}[1]{\textcolor{red}{#1}}
%\newcommand{\green}[1]{\textcolor{green}{#1}}

%\newcommand{\notes}[1]{\marginpar{\textcolor{red}{#1}}}


\author{Ievgen~V.~Bondarenko, Rostyslav~V.~Kravchenko\thanks{The author was partially supported by NSF grants
DMS-0308985, DMS-0456185 and DMS-0600975}}
\title{\textbf{Graph-directed systems and self-similar measures on limit spaces\\ of self-similar groups}}


\sloppy \frenchspacing


\newcommand{\xmo}{X^{-\omega}}
\newcommand{\xo}{X^\omega}
\newcommand{\xz}{X^{\mathbb{Z}}}
\newcommand{\Z}{\mathbb{Z}}
\newcommand{\Rs}{\hat{\mathbb{C}}}
\newcommand{\M}{\mathcal{M}}

\newcommand{\sol}[1][G]{\mathcal{S}_{#1}}
\newcommand{\limGs}[1][G]{\mathcal{X}_{#1}}
\newcommand{\lims}[1][G]{\mathscr{J}_{#1}}
\newcommand{\si}{\mathsf{s}}
\newcommand{\e}{\mathsf{e}}
\newcommand{\muG}{\mu}
\newcommand{\muls}{\mathsf{m}}
\newcommand{\mgr}{meas}

\newcommand{\tile}{\mathscr{T}}
\newcommand{\T}{\mathsf{T}}
\newcommand{\nucl}{\mathcal{N}}
\newcommand{\bim}{\mathscr{W}}
\newcommand{\gr}{\Gamma}
\newcommand{\B}{\mathscr{B}}
\newcommand{\F}{\mathscr{F}}

\newcommand{\resanal}{\mathscr{V}}
\newcommand{\Dom}{\emph{Dom}\,}




\newtheorem{theor}{Theorem}
\newtheorem{prop}[theor]{Proposition}
\newtheorem{cor}[theor]{Corollary}
\newtheorem{lemma}{Lemma}
\theoremstyle{definition}
\newtheorem{defi}{Definition}
\newtheorem{example}{Example}
\newtheorem{remark}{Remark}


\begin{document}
\maketitle

\begin{abstract}
Let $G$ be a group and $\phi:H\rightarrow G$ be a contracting homomorphism from a subgroup $H<G$ of
finite index. V.~Nekrashevych \cite{self_sim_groups} associated with the pair $(G,\phi)$ the limit
dynamical system $(\lims,\si)$ and the limit $G$-space $\limGs$ together with the covering
$\cup_{g\in G}\tile\cdot g$ by the tile $\tile$. We develop the theory of self-similar measures
$\muls$ on these limit spaces. It is shown that $(\lims,\si,\muls)$ is conjugated to the one-sided
Bernoulli shift. Using sofic subshifts we prove that the tile $\tile$ has integer measure and we
give an algorithmic way to compute it. In addition we give an algorithm to find the measure of the
intersection of tiles $\tile\cap (\tile\cdot g)$ for $g\in G$. We present applications to the
invariant measures for the rational functions on the Riemann sphere and to the evaluation of the
Lebesgue measure of integral self-affine tiles.\\

\noindent \textbf{Keywords:} self-similar group, limit space, self-similar measure, tiling,
Bernoulli shift, graph-directed system, self-affine tile.\\


\noindent \textbf{MSC 2000:} Primary 28A80, 20F65; Secondary 37B10, 37D40.\\
\end{abstract}







%Primary 28A80(Fractals), 20F65(GGT); Secondary 37B10(Symbolic dynamics), 37D40 (Dynamical systems
%of geometric origin and hyperbolicity), 37C40 (Smooth ergodic theory, invariant measures).



%A finite state automaton is called deterministic if for every vertex
%the outcomming edges are labeled distinctly.


\section{Introduction}

Let $G$ be a group and $\phi$ be a virtual endomorphism of $G$, which is a homomorphism from a
subgroup $H<G$ of finite index to $G$. Iterative construction involving $\phi$ (together with some
additional data) produces the so-called self-similar action $(G,X^{*})$ of the group $G$ on the
space $X^{*}$ of finite words over an alphabet $X$. And conversely, every self-similar action of
the group $G$ defines a virtual endomorphism of $G$, which almost completely describes the action.
A rich geometric theory is associated with the pair $(G,\phi)$ in \cite{self_sim_groups} through
the theory of self-similar groups. The goal of this paper is to introduce measure to this theory.


Self-similar group is a rather new notion in geometric group theory. Like the self-similar objects
in geometry (fractals) are too irregular to be described using the language of classical Euclidean
geometry, the self-similar groups possess properties not typical for the traditional group theory.
In particular, the class of self-similar groups contains infinite periodic finitely generated
groups, groups of intermediate growth, just-infinite groups, groups of finite width, etc. (see
\cite{self_sim_groups,fractal_gr_sets,branch,gri:MilnorProb} and references therein). At the same
time, it was discovered that self-similar groups appear naturally in many areas of mathematics, and
have applications to holomorphic dynamics, combinatorics, analysis on fractals, etc. An important
class of self-similar groups are contracting groups, which correspond to self-similar actions with
contracting virtual endomorphism. A virtual endomorphism $\phi$ is contracting if it asymptotically
contracts the length of group elements with respect to some generating set. The contracting
property makes many problems around the group effectively solvable.





%A virtual endomorphism is contracting if it asymptotically contracts the length of group elements;
%in this case the associated self-similar actions are also called contracting. Namely the strong
%contracting properties of Grigorchuk group was used to prove that it has intermediate growth. The
%study of the self-similar actions of a group is in some sense the study of iterations of the
%virtual endomorphisms of this group.

%the coset tree $\{g \phi^{-n}(G) | g\in G, n\geq 0\}$





V.~Nekrashevych in \cite{self_sim_groups} associated a limit dynamical system $(\lims,\si)$ with
every contracting self-similar action, where $\lims$ is a compact metrizable space and
$\si:\lims\rightarrow\lims$ is an expanding continuous map. The limit space $\lims$ can be defined
as the quotient of the space of left-infinite sequences $\xmo=\{\ldots x_2x_1 | x_i\in X\}$ by the
equivalence relation, which can be recovered from a finite directed labeled graph $\nucl$, called
the nucleus of the action. Another associated geometric object is the limit $G$-space $\limGs$,
which is a metrizable locally compact topological space with a proper co-compact (right) action of
$G$. The limit spaces $\lims$ and $\limGs$ depend up to homeomorphism only on the pair $(G,\phi)$.
However every self-similar action with the pair $(G,\phi)$ additionally produces a tile $\tile$ of
the limit $G$-space and a covering $\limGs=\cup_{g\in G}\tile\cdot g$ (not a tiling in general).



Limit spaces connect self-similar groups with the classical self-similar sets. The self-similar set
(the attractor) given by the system of contracting similarities $f_1,\ldots,f_n$ (iterated function
system) of a complete metric space is the unique compact set $T$ satisfying $T=\cup_{i=1}^n
f_i(T)$. Given a probability vector $p=(p_1,\ldots,p_n)$, Hutchinson
\cite{hutch:fract_self_similarity} showed the existence of a unique probability measure $\mu$
supported on $T$ satisfying
\begin{equation}\label{eqn_self_similar_measure}
\mu(A)=\sum_{i=1}^{n} p_i\mu(f_i^{-1}(A)), \mbox{ for any Borel set $A$},
\end{equation}
which is called the self-similar measure. Another way to introduce this measure is to consider the
natural coding map $\pi:\xmo\rightarrow T$ given by $\pi(\ldots x_2x_1)=\cap_{m\geq 1} f_{x_1}\circ
f_{x_2}\circ\ldots\circ f_{x_m}(T)$. Then the self-similar measure $\mu$ is the image of the
Bernoulli measure $\mu_p$ on $\xmo$ with weight $p$ (here $\mu_p(x_i)=p_i$) under the projection
$\pi$. Self-similar measures play an important role in the development of fractal geometry, and
have applications in harmonic analysis, conformal dynamics, algebraic number theory, etc. (see
\cite{edgar:int_prob_fr_mes,urb:meas_dim_confdyn, ifs_overl_self_sim_mes,str:self_harm,
bandt:self_sim_meas} and references therein).

%mattila:geometry

The Bernoulli measure $\mu_p$ is one of the basic self-similar measures on the spaces of (left or
right) infinite sequences over the alphabet $X$. In Section~\ref{section_Sofic} we study the
Bernoulli measure of sofic subshifts and other sets given by a finite directed graph $\gr=(V,E)$,
whose edges are labeled by elements of $X$. Consider the set $F_v$ for $v\in V$ of all sequences
$\ldots x_2x_1$, which are read along left-infinite paths ending in the vertex $v$. It is proved in
Section~\ref{section_Sofic} that if the graph $\gr$ is right-resolving (i.e. for every vertex $v$
the outgoing edges at $v$ are labeled distinctly) then the sum $\mgr(\gr)=\sum_{v\in V} \mu_p(F_v)$
is integer, which does not depend on the probability vector $p$. It can be interpreted as follows:
almost every left-infinite sequence belongs to precisely $\mgr(\gr)$ sets $F_v$. The number
$\mgr(\gr)$ we call the measure number of the graph $\gr$. We propose an algorithmic method to
compute the measures $\mu_p(F_v)$ for any graph $\gr$, in particular its measure number.


The push-forward of the uniform Bernoulli measure on $\xmo$ provides the self-similar measure
$\muls$ on the limit space $\lims$. The $G$-invariant measure $\mu$ on the limit $G$-space $\limGs$
is defined in a similar way. The measure $\mu$ restricted to the tile $\tile$ satisfies the
self-similarity equation (\ref{eqn_self_similar_measure}), so it is also self-similar. It is proved
in Section~\ref{section_Measures} that the measure $\mu(\tile)$ is equal to the measure number of
the nucleus $\nucl$. In particular it is integer, the fact which generalizes corresponding result
for integral self-affine tiles \cite{lag_wang:int_self_affine_I}. In addition we give an algorithm
to find the measure of intersection of tiles $\tile\cap (\tile\cdot g)$ for $g\in G$. Then the
covering $\limGs=\cup_{g\in G}\tile\cdot g$ is a perfect multiple covering of multiplicity
$\mu(\tile)$, i.e. every point of $\limGs$ belongs to at least $\mu(\tile)$ tiles and almost every
point belongs to precisely $\mu(\tile)$ tiles. This is used to prove that the measures $\muls$ and
$\mu$ depend not on the specific self-similar action of $G$, but only on the pair $(G,\phi)$ as the
limit spaces themselves. Using a criterion from \cite{hoffman:Bernoulli} we show that the limit
dynamical system $(\lims,\si,\muls)$ is conjugated to the one-sided Bernoulli shift.

This work is partially motivated by applications presented in Section~\ref{section_Applictions}. If
$G$ is a torsion-free nilpotent group with a contractive surjective virtual endomorphism $\phi$ and
a faithful self-similar action, then the measure $\mu$ on $\limGs$ can be considered as a Haar
measure on the respective nilpotent Lie group, Malcev's completion of $G$. In the case of
self-similar actions of the free abelian group $\mathbb{Z}^n$ the limit $G$-space
$\limGs[\mathbb{Z}^n]$ is $\mathbb{R}^n$ and the tile $\tile$ is an integral self-affine tile,
which are intensively studied for the last two decades (see \cite{lag_wang:self_affine,
vince:digit, lag_wang:int_self_affine_I,lag_wang:int_self_affine_II,he_hui:self_affine}). In this
case the measure $\mu$ is the Lebesgue measure on $\mathbb{R}^n$. One can apply the methods
developed in Section~\ref{section_Sofic} to give an algorithmic way to find the Lebesgue measure of
an integral self-affine tile, providing answer to the question in \cite{lag_wang:int_self_affine_I}
(initially solved in \cite{gab_yu:natur_til} without self-similar groups). In addition we have an
algorithm to find the Lebesgue measure of the intersection of tiles $\tile\cap (\tile+a)$ for
$a\in\mathbb{Z}^n$ studied in \cite{gab_yu:natur_til,int_two_tiles}. Finally, the theory of
iterated monodromy groups developed in \cite{self_sim_groups} implies that if $G$ is the iterated
monodromy group of a sub-hyperbolic rational function $f$ then the measure $\muls$ on the limit
space $\lims$ is the unique $f$-invariant probability measure $\mu_f$ on the Julia set of $f$
studied in \cite{invmeas_rat,lub:entr_ration,rational_bernoulli}.


\section{Bernoulli measure of sofic subshifts}\label{section_Sofic}

Let $X$ be a finite set with discrete topology. Denote by $X^{*}=\{x_1x_2\ldots x_n | x_i\in X,
n\geq 0\}$ the set of all finite words over $X$ (including the empty word denoted $\emptyset$). Let
$X^{\omega}$ be the set of all right-infinite sequences (words) $x_1x_2\ldots$, $x_i\in X$. Let
$X^{-\omega}$ be the set of all left-infinite sequences (words) $\ldots x_2x_1$, $x_i\in X$. We put
the product topology on these sets. The length of a word $v=x_1x_2\ldots x_n$ is denoted by $|v| =
n$.

The \textit{shift on the space $\xo$} is the map $\sigma:\xo\rightarrow\xo$, which deletes the
first letter of a word, i.e. $\sigma(x_1x_2x_3\ldots) = x_2x_3\ldots$. The \textit{shift on the
space $\xmo$} is the map also denoted by $\sigma$, which deletes the last letter of a word, i.e.
$\sigma(\ldots x_3x_2x_1) = \ldots x_3x_2$. The shifts are continuous $|X|$\,-\,to\,-$1$ maps. The
branches $\sigma_x$ for $x\in X$ of the inverse of $\sigma$ are defined by
$\sigma_x(x_1x_2\ldots)=xx_1x_2\ldots$ and $\sigma_x(\ldots x_2x_1)=\ldots x_2x_1x$.

Let $p=(p_x)_{x\in X}$ be a probability vector (fixed for the rest of the section) and let $\mu_p$
be the \textit{Bernoulli measure} on $\xo$ with weight $p$, i.e. this measure is defined on
cylindrical sets by
\[
\mu_p(x_1x_2\ldots x_n\xo)=p_{x_1}p_{x_2}\ldots p_{x_n}.
\]
The measure on $\xmo$ is defined in the same way. We always suppose that $p_x>0$ for all $x\in X$
(otherwise we can pass to a smaller alphabet $X$). In case $p_x=\frac{1}{|X|}$ for all $x\in X$,
the measure $\mu_p$ is the \textit{uniform Bernoulli measure} denoted $\mu_u$. The dynamical system
$(\xo,\sigma,\mu_u)$ is called the \textit{one-sided Bernoulli $|X|$-shift}. The measure $\mu_p$ is
the unique regular Borel probability measure on $\xo$ that satisfies the self-similarity condition:
\[
\mu_p(A)=\sum_{x\in X} p_x\mu_p(\sigma_x^{-1}(A))
\]
for any Borel set $A\subset\xo$.

%(further we use notation $p_{x_1x_2\ldots x_n}=p_{x_1}p_{x_2}\ldots p_{x_n}$).

In this section all considered graphs are directed and labeled with $X$ as the set of labels. Let
$\gr=(V,E)$ be such a graph and take a vertex $v\in V$. We say that a sequence $x_1x_2\ldots\in\xo$
(a word $x_1x_2\ldots x_n\in X^{*}$) \textit{starts in the vertex $v$} if there exists a
right-infinite path $e_1e_2\ldots$ (finite path $e_1e_2\ldots e_n$) in $\gr$, which starts in $v$
and is labeled by $x_1x_2\ldots$ (respectively $x_1x_2\ldots x_n$). Similarly, we say that a
sequence $\ldots x_2x_1\in\xmo$ (a word $x_n\ldots x_2x_1\in X^{*}$) \textit{ends in the vertex
$v$} if there exists a left-infinite path $\ldots e_2e_1$ (finite path $e_n\ldots e_2e_1$) in
$\gr$, which ends in $v$ and is labeled by $\ldots x_2x_1$ (respectively $x_n\ldots x_2x_1$). For
every $w\in X^*\cup\xmo$ denote by $V_\gr(w)=V(w)\subset V$ the set of all vertices $v\in V$ such
that the sequence $w$ ends in $v$. Observe that $V(w'w)\subseteq V(w)$ for arbitrary word $w'$ and
finite word~$w$.

For every vertex $v\in V$ denote by $B_v$ the set of all right-infinite sequences that start in
$v$, and denote by $F_v$ the set of all left-infinite sequences that end in $v$. The sets $B_v$ and
$F_v$ are closed correspondingly in $\xo$ and $\xmo$, thus compact and measurable. The sets
$\B=\cup_{v\in V}B_v$ and $\F=\cup_{v\in V}F_v$ are the one-sided (respectively, right and left)
\textit{sofic subshifts} associated with the graph $\gr$. The sets $F_v, v\in V$, satisfy the
recursion
\[
F_v=\bigcup_{u\, \stackrel{x}{\rightarrow} \,v} \sigma_x(F_u)
\]
(here the union is taken over all edges which end in $v$). Hence, associating the map $\sigma_x$
with every edge $e$ of the graph $\gr$ labeled by $x$, the collection of sets $\{F_v, v\in V\}$ can
be seen as the graph-directed iterated function system on the sofic subshift $\F$ with the
underlying graph $\gr$ (see \cite{fractal_gr_sets}). All the maps $\sigma_x^{-1}$  are restrictions
of the shift $\sigma$, and thus $\{F_vx, v\in V\}$ is the Markov partition of the dynamical system
$(\F,\sigma)$. Similarly, the collection of sets $\{B_v, v\in V\}$ can be seen as the
graph-directed iterated function system on the sofic subshift~$\B$.


A labeled graph $\gr=(V,E)$ is called \textit{right-resolving} (Shannon graph in some terminology)
if for every vertex $v\in V$ the edges starting at $v$ have different labels. Every sofic subshift
can be given by a right-resolving graph (see Theorem~3.3.2 in \cite{IntroSymbDyn}). A
right-resolving graph is called \textit{strictly right-resolving} if every vertex $v\in V$ has an
outgoing edge labeled by $x$ for every $x\in X$.



%\begin{proof}
%By changing the direction of every edge of the graph $\gr$ the sets $B_v$ and $F_v$ transform to
%each other. Thus it is sufficient to prove the statement in one direction. Suppose that $\cup_{v\in
%V} B_v=\xo$ and take arbitrary sequence $\ldots x_2x_1\in\xmo$. For each $n\geq 1$ there exists a
%path in $\gr$ labeled by $x_n\ldots x_2x_1$. It was shown in the proof of the theorem that then
%there exists a left-infinite path labeled by $\ldots x_2x_1$. Hence $\F=\xmo$.
%
%
%If the graph $\gr$ is strictly right-resolving then $B_v=\xo$ for every vertex $v$, and thus
%$\F=\xmo$.
%\end{proof}

For a labeled graph $\gr=(V,E)$ we use the following notations:
\[
\vec\mu_p(\B)=(\mu_p(B_v))_{v\in V}, \quad  \vec\mu_p(\F)=(\mu_p(F_v))_{v\in V}, \quad \mbox{ and }
\quad \mu_p(\gr)=\sum_{v\in V} \mu_p(F_v).
\]
Next we study the properties of these quantities and describe an algorithmic way to find them.
First we discuss the problem for right-resolving graphs, and then reduce the general case to the
right-resolving one.




\begin{theor}\label{theor_measure_graph_integer}
Let $\gr=(V,E)$ be a finite right-resolving graph. Then
\[
\mu_p(\gr)=\min_{w\in\xmo} |V(w)|=\min_{w\in X^*} |V(w)|.
\]
In particular, the measure $\mu_p(\gr)$ is integer.
\end{theor}
\begin{proof}
Let $w=\ldots x_2x_1\in\xmo$ and denote $w_n=x_n\ldots x_2x_1$ for $n\geq 1$. Observe, that
$V(w)\subseteq V(w_n)$ and $V(w_n)\subseteq V(w_m)$ for $n\geq m$.

Take a vertex $v\in \cap_{n\geq 1} V(w_n)$. Let $P_n$ be the set of all paths $e_n\ldots e_2e_1$
labeled by $w_n$ and ending in $v$. The set $P_n$ is a finite non-empty set for every $n$, and
$e_{n-1}\ldots e_2e_1\in P_{n-1}$ for every path $e_n\ldots e_2e_1\in P_n$. Since the inverse limit
of a sequence of finite non-empty sets is non-empty, there exists a left-infinite path $\ldots
e_2e_1$ labeled by $w$ and ending in $v$. Then $v\in V(w)$ and we get
\begin{equation}\label{eqn_th1_V(w) eq cap V(w_n)}
V(w)=\bigcap_{n\geq 1} V(w_n).
\end{equation}
From this follows that $|V(w)|=\min\limits_{n\geq 1} |V(w_n)|$ for $w\in\xmo$. Then
\[
\min_{w\in\xmo} |V(w)|=\min_{w\in\xmo} \min_{n\geq 1} |V(w_n)|=\min_{w\in X^*} |V(w)|
\]
and the second equality of the theorem is proved.

Define the integer $k=\min_{w\in\xmo} |V(w)|$ and consider the set
$\mathcal{O}=\mathcal{O}(\gr)\subseteq\xmo$ of all sequences $w\in\xmo$ such that $|V(w)|=k$.
Define $\mathcal{O}^*$ as the set of finite words that satisfy the same condition. We have the
following lemma.
\begin{lemma}\label{lemma_set V open dense}
The set $\mathcal{O}$ is open and dense in $\xmo$, and $\mu_p(\mathcal{O})=1$. For each
$w\in\mathcal{O}$ there is a beginning of $w$ that belongs to $\mathcal{O}^*$. Equivalently,
$\mathcal{O}=\cup_{w\in\mathcal{O}^*}\xmo w$.
\end{lemma}
\begin{proof}
If $w\in\mathcal{O}$ then $k=|V(w)|=\min_{n\geq 1} |V(w_n)|$ and there exists $N\geq 1$ such that
$|V(w_N)|=k$. Then $k\leq |V(\omega w_N)|\leq |V(w_N)|=k$ for all $\omega\in\xmo$. Hence
$w_N\in\mathcal{O}^*$, $\xmo w_N\subseteq\mathcal{O}$ and so $\mathcal{O}$ is open, and thus
measurable.

Let $u\in X^*$ be such that $|V(u)|=k$. Let us show that if $w\in\xmo$ contains the subword $u$
then $w\in\mathcal{O}$. If $w=w'u$ then $k\leq |V(w'u)|\leq |V(u)|=k$ and $w\in\mathcal{O}$.
Observe that $V(ux)$ is the set of those vertices $v\in V$ for which there exists an edge labeled
by $x$ which starts in some vertex of $V(u)$ and ends in $v$. Since the graph $\gr$ is
right-resolving there is no more than one such an edge for each vertex of $V(u)$, and thus
$|V(ux)|\leq |V(u)|$. It implies that if $w=w'uu'$ than $k\leq |V(w)|\leq |V(u)|=k$ and thus
$|V(w)|=k$, so $w\in\mathcal{O}$. The Bernoulli measure of the set of all words $w'uu', u'\in X^*,
w'\in\xmo$, is equal to $1$. Thus $\mu_p(\mathcal{O})=1$. It follows also that $\mathcal{O}$ is
dense in $\xmo$.
\end{proof}


By construction of the set $\mathcal{O}$ for every $w\in\mathcal{O}$ there exist exactly $k$
vertices $v$ such that $w\in F_v$. Let $\chi_{F_v}$ be the characteristic function of the set
$F_v$. Then $\sum_{v\in V}\chi_{F_v}=k$ almost everywhere. Integrating we get
$\mu_p(\gr)=\sum\limits_{v\in V} \mu_p(F_v)=k$.
\end{proof}



%\begin{defi}
%Let $\gr=(V,E)$ be a strictly right-resolving graph. Then we say that $\gr$ satisfies the open set
%condition if there is $v_0\in V$ such that for each $v\in V$ there is a path in $\gr$ from $v$ to
%$v_0$ and for each $x\in X$ there is an edge from $v_0$ to $v_0$ labeled by $x$.
%\end{defi}
%
%\begin{cor}\label{ops}
%If $\gr$ satisfies open set condition, then $\mathcal{V}(\gr)=\{w\in\xmo | \#V_\gr(w)=1\}$ and $\mu_p(\gr)=1$.
%\end{cor}
%\begin{proof}
%Suppose that the action satisfies the open set condition. Let $\gr=\{v_0,\ldots,v_n\}$.
%Choose a word $u_1\in X^{*}$ such that the path labeled by $u_1$ goes from $v_1$ to $v_0$, and for every $v_i$ find consequently a word
%$u_i\in X^{*}$ such that the path labeled by $u_1\ldots u_i$ goes from $v_i$ to $v_0$. Then for the word $u=u_1\ldots u_n$ the
%path labeled by $u$ goes from $v$ to $v_0$ for all $v\in V$. Thus $V(u)=\{v_0\}$. It follows that $\min \#V(w)=1$.
%\end{proof}


\begin{remark} The theorem holds not only for Bernoulli measures. The only property that was used is that
for every word $u\in X^{*}$ the set of all words, which contain $u$ as a subword, has measure $1$.
It should be pointed out that the number $\mu(\gr)$ is independent on the chosen measure $\mu$,
while the measures $\mu(F_v)$, $\mu(B_v)$ and $\sum_{v\in V} \mu(B_v)$ depend on $\mu$.
\end{remark}

\begin{defi}
The number $\mgr(\gr)=\mu_p(\gr)$ is called the \textit{measure number} of the graph~$\gr$.
\end{defi}

Theorem~\ref{theor_measure_graph_integer} shows that almost every sequence $w\in\xmo$ ends in
precisely $\mgr(\gr)$ vertices of $\gr$.


We will use the following proposition in the next sections.
\begin{prop}\label{prop cup B_v eq cup F_v}
Let $\gr=(V,E)$ be a finite labeled graph. Then $\B=\xo$ if and only if $\F=\xmo$. In particular,
$\F=\xmo$ for a finite strictly right-resolving graph $\gr$.
\end{prop}
\begin{proof}
Since the inverse limit of nonempty finite sets is nonempty, $\B=\xo$ ($\F=\xmo$) is equivalent to
the fact that every finite word $v\in{X^*}$ labels some path in $\gr$.
\end{proof}



The matrix $A=(a_{vu})_{v,u\in V}$, where $a_{vu}$ is equal to the number of edges from $v$ to $u$,
is the \textit{adjacency matrix} of the graph $\gr$. For the probability vector $p=(p_x)_{x\in X}$
define the matrix
\[
T_p=(t_{vu})_{v,u\in V}, \mbox{ where } t_{vu}=\sum\limits_{v\stackrel{x}\rightarrow u}p_x
\]
(the sum is taken over all edges from $v$ to $u$). The matrices $A$ and $T_p$ are irreducible if
and only if the graph $\gr$ is strongly connected. If the graph $\gr$ is right-resolving, then the
matrix $T_p$ is the \textit{transition matrix} of the random walk on the weighted directed graph
$\gr$, where each edge labeled by $x$ has weight $p_x$. In this case the row sums of the matrix $A$
are $\leq |X|$, and the row sums of the matrix $T_p$ are $\leq 1$, hence the spectral radius of $A$
is $\leq |X|$, and the spectral radius of $T_p$ is $\leq 1$. If the graph $\gr$ is strictly
right-resolving, then the transition matrix $T_p$ is right stochastic.


\begin{prop}\label{prop_eigenvect_transition_matr}
Let $\gr=(V,E)$ be a finite right-resolving graph with the transition matrix~$T_p$. If the vector
$\vec \mu_p(\B)$ is nonzero then it is the right eigenvector of $T_p$ for the eigenvalue~$1$. If
the vector $\vec \mu_p(\F)$ is nonzero then it is the left eigenvector of $T_p$ for the
eigenvalue~$1$.
\end{prop}
\begin{proof}
By construction, for every vertex $v\in V$ we have
\[
B_v=\bigsqcup_{v\stackrel{x}\rightarrow u} xB_u
\]
(here the union is disjoint because the graph $\gr$ is right-resolving). It implies
\[
\mu_p(B_v)=\sum_{v\stackrel{x}\rightarrow u} p_x\mu_p(B_{u})=\sum_{u\in V} t_{vu}\mu_p(B_{u}).
\]
Thus the nonzero vector $\vec \mu_p(\B)$ is the right eigenvector of $T_p$ for the eigenvalue~$1$.

Similarly, %for every vertex $v\in V$ we have
\[
F_v=\bigcup_{u\stackrel{x}\rightarrow v} F_{u}x,\ v\in V,
\]
and, since the graph $\gr$ is right-resolving, that implies
\begin{equation*}%\label{eqn_sum_prop_eigenvector_of_right_res}
\mu_p(F_v)\leq\sum_{u\stackrel{x}\rightarrow v} p_x\mu_p(F_{u})=\sum_{u\in V}
t_{uv}\mu_p(F_{u})\quad \Rightarrow\quad \vec \mu_p(\F)\leq\vec \mu_p(\F) T_p.
\end{equation*}

The standard arguments based on the theory of nonnegative matrices (see for example \cite[proof of
Theorem 4.5]{he_hui:self_affine}, \cite[p. 197]{gab_yu:natur_til}, \cite{handbook:rothblum}) end
the proof.
%
%
%Let $\gr^{(1)}$ be a strongly connected component of $\gr$. Suppose that there is an outgoing edge
%from $\gr^{(1)}$, and let us show that $\mu_p(F_v)=0$ for every vertex $v$ of $\gr^{(1)}$. By
%induction on the strongly connected components we can suppose that the incoming edges do not carry
%additional measures, or, equivalently, there are no incoming edges to the component $\gr^{(1)}$.
%The transition matrix $T_p^{(1)}$ of $\gr^{(1)}$ is irreducible. Since some row sum of $T_p^{(1)}$
%is less than $1$, the spectral radius of $T_p^{(1)}$ is less than $1$ by Fact~9a in
%\cite[Section~9.3]{handbook:rothblum}. As previously, $\vec \mu_p(\F^{(1)})\leq\vec \mu_p(\F^{(1)})
%T_p^{(1)}$, where the vector $\vec \mu_p(\F^{(1)})$ is the restriction of $\vec \mu_p(\F)$ to the
%vertices of $\gr^{(1)}$. Hence, $\vec \mu_p(\F^{(1)})=0$ by Fact~8d in
%\cite[Section~9.3]{handbook:rothblum} and the statement is proved.
%
%Hence, for every strictly connected component $\gr^{(1)}$ the incoming edges do not carry
%additional measure in the summation~(\ref{eqn_sum_prop_eigenvector_of_right_res}), which implies
%$\vec \mu_p(\F^{(1)})\leq\vec \mu_p(\F^{(1)}) T_p^{(1)}$. Now if the vector $\vec \mu_p(\F^{(1)})$
%is nonzero, then the spectral radius of $T_p^{(1)}$ is equal to $1$, and the vector $\vec
%\mu_p(\F^{(1)})$ is a left eigenvector of $T_p^{(1)}$ for the eigenvalue $1$ by Fact~5d in
%\cite[Section~9.3]{handbook:rothblum}.
%
%It implies that if the vector $\vec \mu_p(\F)$ is nonzero, then the spectral radius of $T_p$ is
%equal to $1$, and $\vec \mu_p(\F)$ is a left eigenvector of $T_p$ for the eigenvalue $1$.
\end{proof}

\begin{cor}\label{cor_meas_sum_meas}
Let $\gr=(V,E)$ be a finite right-resolving graph. Let $\{\gr_i\}$ be the set of all strongly
connected components of $\gr$, which are strictly right-resolving graphs. Then
$\mgr(\gr)=\sum_i\mgr(\gr_i)$.
\end{cor}

In particular, if a finite strictly right-resolving graph $\gr$ contains a vertex $v_0$ such that
for each vertex $v$ there is a path in $\gr$ from $v$ to $v_0$ and for each $x\in X$ there is an
edge from $v_0$ to $v_0$ labeled by $x$ (the open set condition for graphs), then $\mgr(\gr)=1$.


\begin{cor}
Let $\gr=(V,E)$ be a finite right-resolving graph with the adjacency matrix~$A$. For the uniform
Bernoulli measure $\mu_u$ the nonzero vectors $\vec \mu_u(\B)$ and $\vec \mu_u(\F)$ are
respectively the right and left eigenvectors of $A$ for the eigenvalue $|X|$.
\end{cor}


The next example shows that the right-resolving property in
Proposition~\ref{prop_eigenvect_transition_matr} and Theorem~\ref{theor_measure_graph_integer}
cannot be dropped even in the case of the uniform Bernoulli measure. The second example presents
two right-resolving graphs with the same transition matrices, but with different measure numbers.
\begin{example}
For the graph $\gr$ shown in Figure~\ref{fig_Example_3_Graphs} we get $F_a=\xmo 0$, $F_b=\xmo$, the
vector $\vec\mu_u(\F)=(1/2,1)$ is not an eigenvector of the adjacency matrix, and $\mu_u(\gr)=3/2$.
\end{example}

\begin{example}
Consider two graphs $\gr_1$ and $\gr_2$ shown in Figure~\ref{fig_Example_3_Graphs}. They have the
same transition matrix (for the uniform Bernoulli measure), however $F_{a_1}=F_{b_1}=\xmo$ and
$F_{a_2}=\xmo 0$, $F_{b_2}=\xmo 1$. Hence, $\mu_u(\gr_1)=2$ and $\mu_u(\gr_2)=1$.
\end{example}

\begin{figure}
\begin{center}
\psfrag{0}{\small{$0$}} \psfrag{1}{\small{$1$}} \psfrag{a}{\small{$a$}} \psfrag{b}{\small{$b$}}
\psfrag{e}{\small{$a_1$}} \psfrag{f}{\small{$b_1$}} \psfrag{c}{\small{$a_2$}}
\psfrag{d}{\small{$b_2$}} \epsfig{file=fig1.eps,height=70pt}\caption{The graphs $\gr$ (on the
left), $\gr_1$ (in the center), and $\gr_2$ (on the right)} \label{fig_Example_3_Graphs}
\end{center}
\end{figure}


Although Theorem~\ref{theor_measure_graph_integer} gives a usefull characterization of the number
$\mu_p(\gr)$, it does not present an algorithmic way to find it. It follows from
Proposition~\ref{prop_eigenvect_transition_matr} that the problem of finding $\vec\mu_p(\F)$ and
$\mu_p(\gr)$ reduces to the strongly connected components which are strictly right-resolving graphs
(for all other vertices $\mu_p(F_v)=0$). Notice that if $\gr$ is a strongly connected strictly
right-resolving graph, then the vector $\frac{1}{\mu_p(\gr)}\vec \mu_p(\F)$ is the unique
stationary probability distribution of the stochastic matrix $T_p$.

%Hence the problem of finding $\vec\mu_p(\F)$ is equivalent to the problem of finding $\mu_p(\gr)$.

At the same time Proposition~\ref{prop_eigenvect_transition_matr} implies the algorithm to find the
vector $\vec\mu_p(\B)$ for a right-resolving graph. Indeed, a left eigenvector of $T_p$ for the
eigenvalue $1$ is uniquely defined if we know its entries $\mu_p(B_v)$ for vertices $v$ in the
strongly connected components of $\gr$ without outgoing edges. For every such a component $\gr'$,
we have $B_v=\xo$ and $\mu_p(B_v)=1$ for every vertex $v\in\gr'$ if the component $\gr'$ is a
strictly right-resolving graph, and $\mu_p(B_v)=0$ otherwise. In particular, if the matrix $T_p$ is
rational then the values $\mu_p(B_v)$ are rational.

%However, since the transition (adjacency) matrix does not capture the edge labeling, it is not
%enough to find $\mu_p(\gr)$ as the following example shows.




%\begin{prop}\label{prop_mu_lim_Tp}
%Let $\gr=(V,E)$ be a finite right-resolving graph with the transition matrix~$T_p$. Then
%\[
%\vec \mu_p(\B)=\lim_{n\rightarrow\infty} T_p^n\vec {1},
%\]
%where $\vec {1}=(1,1,\ldots,1)$.
%\end{prop}
%\begin{proof}
%Consider the approximations $B_v^{(n)}$ of the sets $B_v$ given by the graph-directed iterated
%function construction: put $B_v^{(0)}=\xo$ for every $v\in V$ and define inductively
%\[
%B_v^{(n+1)}=\bigsqcup_{v\stackrel{x}\rightarrow u} xB_u^{(n)}\quad \Rightarrow\quad
%\mu_p(B_v^{(n+1)})=\sum_{v\stackrel{x}\rightarrow u} p_{x}\mu_p(B_u^{(n)})=\sum_{u\in V}
%t_{vu}\mu_p(B_u^{(n)})
%\]
%(here the union is disjoint because the graph $\gr$ is right-resolving). Hence the vector $\vec
%\mu_p(B^{(n)})=(\mu_p(B_v^{(n)}))_{v\in V}$ can be found by
%\[
%\vec \mu_p(B^{(n)})=T_p\vec \mu_p(B^{(n-1)}) =\ldots =T_p^n\vec \mu_p(B^{(0)})=T_p^n\vec {1}.
%\]
%The sets $B_v^{(n)}$ form a decreasing sequence of sets of finite measure and $B_v=\cap_{n\geq
%1}B_v^{(n)}$ (the proof is the same as of (\ref{eqn_th1_V(w) eq cap V(w_n)})). Hence
%$\mu_p(B_v)=\lim_{n\rightarrow\infty}\mu_p(B_v^{(n)})$.
%\end{proof}
%
%
%
%\begin{remark}
%Notice that the limit can be found algorithmically. For example, we can diagonalize the matrix and
%find the limit.
%\end{remark}




Given a finite labeled graph $\gr=(V,E)$ the problems of finding the measures $\mu_p(B_v)$ and
$\mu_p(F_v)$ are equivalent, and can be reduced to right-resolving graphs. The problem is that the
system
\[
B_v=\bigcup_{v\stackrel{x}\rightarrow u} xB_{u}=\bigsqcup_{x\in
X}x\left(\bigcup_{v\stackrel{x}\rightarrow u} B_{u}\right),\ v\in V
\]
is not self-similar in the sense that the above expression involves not only the sets $B_v$ but
also their finite unions. Introducing additional terms to this recursion corresponding to these
unions we get a system with a right-resolving graph. This procedure is similar to the construction
described in the proof of Theorem~3.3.2 in \cite{IntroSymbDyn}.


\begin{prop}
For every finite graph $\gr=(V,E)$ one can construct a finite right-resolving graph $\gr'=(V',E')$
with the property that for every $v\in V$ there exists $v'\in V'$ such that $B_v=B_{v'}$.
\end{prop}
%\begin{proof}
%In the new graph all vertices are subsets of $V$ (one-element subsets are identified with elements
%of $V$). Starting from the graph $\gr$, we can proceed inductively as follows: for every $x\in X$
%and every vertex $v$ of already constructed graph:
%\begin{enumerate}
%  \item Take the set $v'=\{\mbox{elements of } u |v\stackrel{x}\rightarrow u\}\subset V$.
%  If this set was not constructed before, consider $v'$ as a new vertex.
%  \item Remove all the edges $v\stackrel{x}\rightarrow u$ and add one edge $v\stackrel{x}\rightarrow
%  v'$.
%  \item Add edges $v'\stackrel{y}\rightarrow
%w$ for every $v\stackrel{x}\rightarrow u\stackrel{y}\rightarrow w$, $y\in X$.
%\end{enumerate}
%Under this construction, the sets $B_v$, $v\in V$, remain the same, and $B_{v'}=\bigcup_{v\in
%v'}B_v$. Since the graph $\gr$ is finite, after finitely many steps we get a right-resolving graph
%$\gr'=(V',E')$, which produces the same sets $B_v$ for $v\in V\subset V'$. We get the following
%statement.
%
%\end{proof}


Similarly, one can find the measures of subshifts $\B$ and $\F$ by introducing new vertices
corresponding to $\cup_{v\in V} B_v$ and $\cup_{v\in V} F_v$.

%\begin{example}
%Give example which will be used in application to self-affine tiles.
%\end{example}

\begin{cor}\label{cor_meas_gr_rational}
Let $\gr=(V,E)$ be a finite labeled graph. For the uniform Bernoulli measure $\mu_u$ all measures
$\mu_u(B_v)$, $\mu_u(F_v)$, $\mu_u(\gr)$ are rational.
\end{cor}
%\begin{proof}
%As shown above, the general case reduces to right-resolving graphs. It follows from the algor If
%the graph $\gr$ is right-resolving, then $\mu_u(\gr)$ is integer by
%Theorem~\ref{theor_measure_graph_integer} and $\frac{1}{\mu_u(\gr)}\vec \mu_{u}(\F)$ is rational as
%the stationary probability vector of the matrix $T_u$ with rational entries.
%\end{proof}


Consider the question how to find the measure of the intersection $B_v\cap B_u$ for $v,u\in V$.
Construct a new graph $\mathscr{G}$ with the set of vertices $V\times V$ and put an edge from
$(v,u)$ to $(v',u')$ labeled by $x\in X$ for every edges $v\stackrel{x}\rightarrow v'$ and
$u\stackrel{x}\rightarrow u'$ in the graph $\gr$ (label products of graphs by Definition~3.4.8 in
\cite{IntroSymbDyn}). It is easy to see that then $B_{(v,u)}=B_v\cap B_u$ (see Proposition~3.4.10
in \cite{IntroSymbDyn}).


The described methods work in a more general case when labels are words over $X$. Let $\gr=(V,E)$
be a directed graph labeled by labels from $X^{*}\setminus\{\emptyset\}$. Define the sets $B_v$ and
$F_v$ for $v\in V$ as above. Similarly, the collection of sets $\{B_v, v\in V\}$ can be seen as the
graph-directed function system with the underlying graph $\gr$, where the map $\sigma_{x_1\ldots
x_n}=\sigma_{x_n}\circ\ldots\circ\sigma_{x_1}$ is associated with every edge labeled by $x_1\ldots
x_n$. To find the measures of the sets $B_v$ and $F_v$, we can replace each edge labeled by a word
$x_1x_2\ldots x_k$ by a sequence of edges labeled by letters $x_1,x_2,\ldots,x_k$, and apply the
above method to this new graph. Also we can directly use
Proposition~\ref{prop_eigenvect_transition_matr}, where such graph $\gr$ is called
\textit{right-resolving} if for every vertex $v\in V$ the word-label of an edge starting at $v$ is
not a prefix of the word-label of another edge starting at $v$ (i.e. in some sense the associated
automaton is deterministic). The transition matrix $T_p$ is defined similarly
\[
T_p=(t_{vu})_{v,u\in V}, \mbox{ where } t_{vu}=\sum\limits_{v\xrightarrow{x_1...x_k}
u}p_{x_1}p_{x_2}\ldots p_{x_k}.
\]
%(here $p$ is the probability vector).






%----------------------------------------------------------------------------------
%----------------------------------------------------------------------------------
%----------------------------------------------------------------------------------


\section{Self-similar actions and its limit spaces}\label{section_Self-similar_groups}
We review in this section the basic definitions and theorems concerning self-similar groups. For a
more detailed account and for the references, see~\cite{self_sim_groups}.

\textbf{Self-similar actions.} A faithful action of a group $G$ on the set $X^*$ is called
\emph{self-similar} if for every $g\in G$ and $x\in X$ there exist $h\in G$ and $y\in X$ such that
\[
g(xw)=yh(w)
\]
for all $w\in X^*$. The element $h$ is called the \textit{restriction} of $g$ on $x$ and denoted
$h=g|_x$. Inductively one defines the restriction $g|_{x_1x_2\ldots x_n}=g|_{x_1}|_{x_2}\ldots
|_{x_n}$ for every word $x_1\ldots x_n\in X^n$. Notice that $(g\cdot h)|_v=g|_{h(v)}\cdot h|_v$ (we
are using left actions).


%With every self-similar action $(G,X^{*})$ is associated a
%\textit{virtual endomorphism}, which is a homomorphism from a
%subgroup of finite index.


\textbf{Virtual endomorphisms.} The study of the self-similar actions of a group is in some sense
the study of the \textit{virtual endomorphisms} of this group, which are homomorphisms from a
subgroup of finite index to the group. There is a general way to construct a self-similar
representation of a group with a given associated virtual endomorphism. Let $\phi: H\rightarrow G$
be a virtual endomorphism of the group $G$, where $H<G$ is a subgroup of index $d$. Let us choose a
left coset transversal $T=\{g_0,g_1,\ldots,g_{d-1}\}$ for the subgroup $H$, and a sequence
$C=\{h_0,h_1,\ldots,h_{d-1}\}$ of elements of $G$ called a \emph{cocycle}. The self-similar action
$(G,X^{*})$ with the alphabet $X=\{x_0,x_1,\ldots,x_{d-1}\}$ defined by the triple $(\phi,T,C)$ is
given by
\[
g(x_i)=x_j,\qquad g|_{x_i}=h_j^{-1}\phi(g_j^{-1}gg_i)h_i,
\]
where $j$ is such that $g_j^{-1}gg_i\in H$ (such $j$ is unique). The action may be not faithful,
the kernel can be described using Proposition~2.7.5 in \cite{self_sim_groups}.


Conversely, every self-similar action can be obtained in this way. Let $(G,X^{*})$ be a
self-similar action and take a letter $x\in X$. The stabilizer $St_G(x)$ of the letter $x$ in the
group $G$ is a subgroup of index $\leq |X|$ in $G$. Then the map $\phi_x: g\mapsto g|_x$ is a
homomorphism from $St_G(x)$ to $G$ called \emph{the virtual endomorphism associated to the
self-similar action}. Choose $T=\{g_y : y\in X\}$ and $C=\{h_y : y\in X\}$ such that $g_y (x)=y$
and $h_y=\left(g_y|_x\right)^{-1}$. Then $T$ is a coset transversal for the subgroup $St_G(x)$ and
the self-similar action $(G,X^*)$ is defined by the triple $(\phi_x,T,C)$. Different self-similar
actions of the group $G$ with the same associated virtual endomorphism are conjugated by
Proposition~2.3.4 in \cite{self_sim_groups}.

%Notice that here the virtual endomorphism associated with a letter $x\in X$ is given by
%$\phi_x(g)=h_x\phi(g_x^{-1}gg_x)h_x^{-1}$.

%\begin{prop}
%The self-similar actions $(G,X^*)$ and $(G,Y^*)$ defined
%respectively by the pairs $(T,C)$ and $(T',C')$ are conjugate and
%the conjugating isomorphism is the map $\alpha : X^*
%\longrightarrow Y^*$ such that
%\[
%\alpha\cdot x=y\cdot s_x\alpha,
%\]
%where $s_x=h'_y\phi({g'}_y^{-1}g_x)h_x^{-1}$ and y is such that
%$g_x^{-1}g'_y\in\Dom\phi$.
%\end{prop}


\textbf{Contracting self-similar actions.} An important class of self-similar actions are
contracting actions. A self-similar action of a group $G$ is called \emph{contracting} if there
exists a finite set $\nucl$ such that for every $g\in G$ there exists $k\in\mathbb{N}$ such that
$g|_v\in\nucl$ for all words $v\in X^*$ of length $\geq k$. The smallest set $\nucl$ with this
property is called the \emph{nucleus} of the self-similar action. The nucleus itself is
self-similar in the sense that $g|_v\in\nucl$ for every $g\in\nucl$ and $v\in X^{*}$. It can be
represented by the \textit{Moore diagram}, which is the directed labeled graph with the set of
vertices $\nucl$, where there is an edge from $g$ to $g|_x$ labeled $(x,g(x))$ for every $x\in X$
and $g\in\nucl$. We identify the nucleus with its Moore diagram, also denoted by $\nucl$. The
contracting property of the action depends only on the virtual endomorphism but not on the chosen
coset transversal and cocycle (see Corollary~2.11.7 in \cite{self_sim_groups}). Notice that every
contracting self-similar group is countable.

Self-similar groups are related to self-similar sets through the notion of limit spaces.

\textbf{Limit $G$-spaces.} Let us fix a contracting self-similar action $(G,X^{*})$. Consider the
space $\xmo\times G$ of all sequences \ $\ldots x_2x_1\cdot g$, $x_i\in X$ and $g\in G$, with the
product topology of discrete sets $X$ and $G$. Two elements $\ldots x_2x_1\cdot g$ and $\ldots
y_2y_1\cdot h$ of $\xmo\times G$ are called \textit{asymptotically equivalent} if there exist a
finite set $K\subset G$ and  a sequence $g_n\in K, n\geq 1,$ such that
\begin{displaymath}
g_n(x_n x_{n-1}\ldots x_1)=y_n y_{n-1}\ldots y_1\qquad \mbox{ and }\qquad g_n|_{x_n x_{n-1}\ldots
x_1}\cdot g=h
\end{displaymath}
for every $n\geq 1$. This equivalence relation can be recovered from the nucleus $\nucl$ of the
action (Proposition~3.2.6 in \cite{self_sim_groups}).


\begin{prop}\label{prop_asympt_Moore_diagr}
Two elements $\ldots x_2x_1\cdot g$ and $\ldots y_2y_1\cdot h$ of $\xmo\times G$ are asymptotically
equivalent if and only if there exists a left-infinite path $\ldots e_2e_1$ in the nucleus $\nucl$
ending in the vertex $hg^{-1}$ such that the edge $e_i$ is labeled by $(x_i,y_i)$.
\end{prop}

The quotient of the set $\xmo\times G$ by the asymptotic equivalence relation is called the
\textit{limit $G$-space} of the action and denoted $\limGs[(G,X^*)]$. The group $G$ naturally acts
on the space $\limGs[(G,X^*)]$ by multiplication from the right.


The map $\tau_x$ defined by the formula
\[
\tau_x(\ldots x_2 x_1\cdot g)=\ldots x_2 x_1 g(x)\cdot g|_x
\]
is a well-defined continuous map on the limit $G$-space $\limGs$ for every $x\in X$, which is not a
homeomorphism in general. Inductively one defines $\tau_{x_1x_2\ldots
x_n}=\tau_{x_n}\circ\tau_{x_{n-1}}\circ\ldots\circ\tau_{x_1}$.

%Observe that $\tau_{y}(t)=\tau_{x}(t\cdot r)\cdot(r|_x)^{-1}$ for any $r\in G$ with $r(x)=y$.

The image of $\xmo\times 1$ in $\limGs$ is called the \textit{(digit) tile} $\tile$ of the action.
The image of $\xmo v\times 1$ for $v\in X^n$ is called the \textit{tile} $\tile_v$, equivalently
$\tile_v=\tau_v(\tile)$. It follows directly from definition that
\[
\limGs=\bigcup_{g\in G} \tile\cdot g \ \ \ \ \mbox{ and } \ \ \ \ \tile=\bigcup_{v\in X^n} \tile_v.
\]
Two tiles $\tile\cdot g$ and $\tile\cdot h$ intersect if and only if $gh^{-1}\in\nucl$. A
contracting action $(G,X^*)$ satisfies the \emph{open set condition} if for any element $g$ of the
nucleus $\nucl$ there exists a word $v\in X^*$ such that $g|_v=1$, i.e. in the nucleus $\nucl$
there is a path from any vertex to the trivial state. If the action satisfies the open set
condition then the tile $\tile$ is the closure of its interior, and any two different tiles have
disjoint interiors; otherwise every tile $\tile\cdot g$ is covered by the other tiles (see
Proposition~3.3.7 in \cite{self_sim_groups}).


%the graph $\nucl_o$ satisfies the open set condition.


The tile $\tile$ and the partition of $\limGs$ on tiles $\tile\cdot g$ depend on the specific
self-similar action of the group $G$. However, up to homeomorphism the limit $G$-space
$\limGs[(G,X^*)]$ is uniquely defined by the associated virtual endomorphism $\phi$ of the group,
hence we denote it by $\limGs[G](\phi)$ (or $\limGs$ for short).

\begin{theor}\label{theor_limsX hom limsY}
Let $\phi: H\rightarrow G$ be a virtual endomorphism of the group $G$. Let $(G,X^*)$ and $(G,Y^*)$
be the contracting self-similar actions defined respectively by the triples $(\phi,T,C)$ and
$(\phi,T',C')$. Then $\limGs[(G,X^*)]$ and $\limGs[(G,Y^*)]$ are homeomorphic and the homeomorphism
is the map $\alpha : \limGs[(G,X^*)] \longrightarrow \limGs[(G,Y^*)]$ such that
\[
\alpha(\tau_x(t))=\tau_y(\alpha(t))\cdot s_x,\ \mbox{ for }\ t\in \limGs[(G,X^*)],
\]
where $s_x={h'}_y^{-1}\phi({g'}_y^{-1}g_x)h_x$ and $y$ is such that $g'_yg_x^{-1}\in H$.
\end{theor}
\begin{proof}
The statement follows from Sections~2.1--2.5 in \cite{self_sim_groups}.
\end{proof}


\textbf{Limit dynamical system.} The factor of the limit $G$-space $\limGs$ by the action of the
group $G$ is called the \textit{limit space} $\lims=\lims(\phi)$. It follows from the definition
that we may also consider $\lims$ as a factor of $\xmo$ by the following equivalence relation: two
left-infinite sequences $\ldots x_2 x_1,\ldots y_2 y_1$ are equivalent if and only if there exists
a left-infinite path $\ldots e_2 e_1$ in the nucleus $\nucl$ such that the edge $e_i$ is labeled by
$(x_i,y_i)$. The limit space $\lims$ is compact, metrizable, finite-dimensional space. It is
connected if the group $G$ is finitely generated and acts transitively on $X^n$ for all $n$.



The equivalence relation on $\xmo$ is invariant under the shift $\sigma$, therefore $\sigma$
induces a continuous surjective map $\si:\lims\rightarrow\lims$, and every point of $\lims$ has at
most $|X|$ preimages under $\si$. The dynamical system $(\lims,\si)$ is called the \textit{limit
dynamical system} of the self-similar action.

The image of $\xmo v$ for $v\in X^{n}$ in $\lims$ is called the \textit{tile} $\T_v$ of the $n$-th
level. Clearly
\[
\T_v=\bigcup_{x\in X}\T_{xv}\qquad\mbox{ and }\qquad \si(\T_{vx})=\T_v
\]
for every $v\in X^{*}$ and $x\in X$. Two tiles $\T_v$ and $\T_u$ of the same level have nonempty
intersection if and only if there exists $h\in\nucl$ such that $h(v)=u$. Under the open set
condition, every tile $\T_v$ is the closure of its interior, and any two different tiles of the
same level have disjoint interiors (Proposition~3.6.5 in \cite{self_sim_groups}). It will be used
in the next section that
\begin{equation}\label{eqn_diam_tile_goes_zero}
\lim_{n\rightarrow\infty} \max_{v\in X^n} \rm{diam} (\T_v)=0
\end{equation}
for any chosen metric on the limit space $\lims$ (see Theorem~3.6.9 in \cite{self_sim_groups}).

The inverse limit of the topological spaces $\lims\stackrel{\si}{\leftarrow} \lims
\stackrel{\si}{\leftarrow} \cdots$ is called the \textit{limit solenoid} $\sol$. One can consider
$\sol$ as a factor of the space $\xz$ of two-sided infinite sequences by the equivalence relation,
where two sequences $\xi,\eta$ are equivalent if and only if there exist a two-sided infinite path
in the nucleus labeled by the pair $(\xi,\eta)$. The two-sided shift on $\xz$ induces a
homeomorphism $\e:\sol\rightarrow\sol$.


Consider a simple example to illustrate the main notions of this section (see
\cite[Chapter~6]{self_sim_groups} for many other examples).

\textbf{Example.} Consider the transformation $a$ of the space $X^{*}$ over the binary alphabet
$X=\{0,1\}$ defined recursively by the rule
\begin{eqnarray*}
a(0w)=1w \quad \mbox{ and }\quad a(1w)=0a(w)
\end{eqnarray*}
for $w\in X^{*}$. The action of $a$ on the space $X^{*}$ corresponds to addition of one in the
binary numeration system, and therefore the transformation $a$ is called the (binary)
\textit{adding machine}. Indeed, $a(x_1x_2\ldots x_n)=y_1y_2\ldots y_n$ for $x_i,y_i\in X$ if and
only if
\[
1+x_1+2x_2+\ldots +2^{n-1}x_n=y_1+2y_2+\ldots+2^{n-1}y_n \mod 2^{n}.
\]
\begin{figure}
\begin{center}
\psfrag{a}{$a$} \psfrag{a1}{$a^{-1}$} \psfrag{id}{$id$} \psfrag{00}{$(0,0)$} \psfrag{11}{$(1,1)$}
\psfrag{01}{$(0,1)$} \psfrag{10}{$(1,0)$} \epsfig{file=fig2.eps,height=110pt}\caption{The nucleus
of the adding machine} \label{fig_AddingMachine}
\end{center}
\end{figure}
The group generated by $a$ is the infinite cyclic group, and we have a self-similar action of the
group $\Z\cong\langle a\rangle$ on the space $X^{*}$. This action corresponds to the virtual
endomorphism $\phi:2\Z\rightarrow\Z$ given by $\phi(m)=m/2$ and the coset transversal $T=\{0,1\}$
(there is no need to specify the cocycle for abelian groups). The nucleus of the action
$(\Z,X^{*})$ is shown in Figure~\ref{fig_AddingMachine}. It follows that two sequences in
$\xmo\times \Z$ are asymptotically equivalent if and only if they are of the form
\[
\ldots 000\cdot a^{n+1} \approx \ldots 111\cdot a^n \quad \mbox{or} \quad \ldots 001w\cdot a^n
\approx \ldots 110w\cdot a^n
\]
for $w\in X^{*}$ and $n\in\Z$. This is the usual identification of binary expansions of real
numbers (but written to the left). Hence the map $\Phi$ given by the rule
\[
\Phi(\ldots x_2x_1\cdot a^n)=n+\sum_{i=1}^{\infty} \frac{x_i}{2^i},\quad x_i\in X \mbox{ and }
n\in\Z,
\]
induces a homeomorphism between the limit $G$-space $\limGs[(\Z,X^{*})]$ and the real line
$\mathbb{R}$ with the standard action of $\Z$ by translations. The image of the tile $\tile$ under
$\Phi$ is the unit interval $[0,1]$. The limit space $\lims[(\Z,X^{*})]$ is homeomorphic to the
unit circle.

%two sequences $\ldots x_2x_1\cdot a^n$ and $\ldots y_2y_1\cdot a^m$ are asymptotically equivalent
%if and only if
%\[
%n+\sum_{i=1}^{\infty} \frac{x_i}{2^i}=m+\sum_{i=1}^{\infty} \frac{y_i}{2^i}.
%\]




%----------------------------------------------------------------------------------
%----------------------------------------------------------------------------------
%----------------------------------------------------------------------------------

\section{Self-similar measures on limit spaces}\label{section_Measures}

Let us fix a contracting self-similar action $(G,X^*)$.\\

\noindent\textbf{Invariant measure on the limit $G$-space $\limGs$}. We consider the uniform
Bernoulli measure $\mu_u$ on the space $\xmo$ and the counting measure on the group $G$, and we put
the product measure on the space $\xmo\times G$. The push-forward of this measure under the factor
map $\pi_{\limGs[]}:\xmo\times G\rightarrow\limGs$ defines the \textit{measure $\mu$ on the limit
$G$-space} $\limGs$. The measure $\mu$ is a $G$-invariant $\sigma$-finite regular Borel measure on
$\limGs$.

\begin{prop}\label{prop_measure lims propert}
The measures of tiles have the following properties.
\begin{enumerate}
  \item $\mu(\tile_v)= |X|^n\cdot\mu(\tile_{uv})$ for every $v\in X^{*}$ and $u\in X^n$.
  \item $\mu(\tile_v\cap\tile_{v'})=0$ for $v,v'\in X^n, v\neq v'$.
  \item Let $\mu|_{\tile}$ be the measure $\mu$ restricted to the tile $\tile$. Then
\[
  \mu|_{\tile}(A)=\sum_{x\in X} \frac{1}{|X|}\mu|_{\tile}(\tau_x^{-1}(A))
\]
for any Borel set $A\subset\tile$.
\end{enumerate}
\end{prop}
\begin{proof}
Let us show that $\mu(A)\geq |X|\mu(\tau_x(A))$ for any Borel set $A$. Consider sequences which
represent points of the sets $A$ and $\tau_x(A)$:
\[
\pi^{-1}_{\limGs[]}(A)=\bigcup_{g\in G}T_g\cdot g \ \ \Rightarrow \ \
\pi^{-1}_{\limGs[]}(\tau_x(A))=\bigcup_{g\in G}T_g g(x)\cdot g|_x.
\]
% Indeed, suppose $\dots{}x_1\cdot{}t\in\pi^{-1}(\tau_x(A))$.  It means that
% exists $\dots{}y_1\cdot{}h\in{}T_h\cdot{}h$ such that $\dots{}y_1{}h(x)h|_x$
% is equivalent to $\dots{}x_1\cdot{}t$.  I.e. There is an el-t $g$ in nucleus
% and an arrow from $g$ to $t(h|_x)^{-1}$ labeled by $(h(x),x_1)$.  Thus
% $g|_{h(x)}=t(h|_x)^{-1}$. It follows that $t=(gh)|_x$ and
% $\dots{}x_3x_2\cdot{gh}$ is equiv-nt to $\dots{}y_2y_1\cdot{h}$, thus
% $\dots{}x_3x_2\cdot{gh}\in{}T_{gh}$. So
% $\pi^{-1}(\tau_x(A))\subset{}\cup_gT_gg(x)\cdot{g|_x}$.
It implies
\[
\mu(\tau_x(A))\leq \sum_{g\in G} \frac{1}{|X|}\mu_u (T_g)=\frac{1}{|X|}\mu(A).
\]

By applying this inequality $n$ times we get
$|X|^n\mu(\tile_v)=|X|^n\mu(\tau_v(\tile))\leq\mu(\tile)$ for $v\in{X}^n$. Since $\tile=\cup_{u\in
X^n}\tile_{u}$ we have that
\[
\mu(\tile)\leq\sum_{u\in X^n}\mu(\tile_{u})\leq\sum_{u\in X^n}\frac{1}{|X|^n}\mu(\tile)=\mu(\tile).
\]
Hence all the above inequalities are actually equalities, $\mu(\tile)=|X|^n\mu(\tile_{u})$ for
every $u\in X^n$, and $\mu(\tile_u\cap\tile_{u'})=0$ for different $u,u'\in X^n$.

Notice that since every Borel set $A\subset\tile$ can be approximated by unions of tiles of the
same level %see proof of theorem~\ref{theor_measure_limGs_invariant}
and using items 1 and 2 we have that if $\mu(\tau_x(A))<\varepsilon$ then $\mu(A)<\varepsilon|X|$.

It is left to prove item 3. First, let us show that the assertion holds for the tiles $\tile_v$.
Since the measure $\mu|_{\tile}$ is concentrated on the tile $\tile$, up to sets of measure zero
the set $\tau_x^{-1}(\tile_v)$ is equal $\tile_u$ if $v=ux$, and is empty if the last letter of $v$
is not $x$. Really, if $t\in\tau_x^{-1}(\tile_{ux})$ and $t\in\tile_v$ with $v\neq u, |v|=|u|$,
then $\tau_x(t)\in\tile_{vx}\cap\tile_{ux}$ and the measure of such points is zero. Hence
\[
\mu|_{\tile}(\tile_v)=\frac{1}{|X|}\mu|_{\tile}(\tile_u)=\sum_{x\in X}
\frac{1}{|X|}\mu|_{\tile}(\tau_x^{-1}(\tile_v)).
\]
Now we can approximate any Borel set by unions of tiles and pass to the limit.
\end{proof}




%\begin{prop}\label{prop_measure lims propert}
% The following properties hold for any
%Borel set $A$, $x\in X$ and $g\in G$:
%\[
%\muG(A)\geq |X|\cdot\muG(\tau_x(A)),
%%\ \ \muG(\tau_x(A))=\muG(\tau_y(A))
%\]
%with the equality when $A=\tile_v\cdot g$, $v\in X^*$. If $v\neq v'\in X^n$ then $\mu(\tile_v\cap\tile_{v'})=0$.
%\end{prop}
%\begin{proof}
%Let us show $\muG(A)\geq |X|\muG(\tau_x(A))$, the other properties are obvious. Consider sequences
%which represent points of the sets $A$ and $\tau_x(A)$:
%\[
%\pi^{-1}(A)=\bigcup_{g\in G}T_g\cdot g \ \ \Rightarrow \ \ \pi^{-1}(\tau_x(A))=\bigcup_{g\in G}T_g
%g(x)\cdot g|_x.
%\]
%It implies
%\[
%\muG(\tau_x(A))\leq \sum_{g\in G} \frac{1}{|X|}\mu_u (T_g)=\frac{1}{|X|}\muG(A).
%\]
%
%Let $v=x_1\dots x_n$ be a word over $X$. Denote by $\tau_v=\tau_{x_1}\circ\dots\circ\tau_{x_n}$.
% By above, $\mu(A)\geq |X|^n\mu(\tau_v(A))$ for any Borel $A$. Since $\tile=\cup_{v\in X^n}\tau_v(\tile)$ we have that
% $$\mu(\tile)\leq\sum_{v\in X^n}\mu(\tau_v(\tile))\leq\sum_v\frac{1}{|X|^n}\mu(\tile)=\mu(\tile)$$
%
%Thus all the above 'less or equal' signs are actually equalities, so $\mu(\tile)=|X|^n\mu(\tau_v(\tile))$ for any $v\in X^*$, and for any $v\neq v'\in X^n$ $\mu(\tile_v\cap\tile_{v'})=0$. It follows, that $\mu(A)=|X|\mu(\tau_x(A))$ for any $A=\tile_v$, where $v\in X^*$ and any $x\in X$.
%
%\end{proof}

The tile $\tile$ of the limit $G$-space $\limGs$ can be considered as the attractor of the iterated
function system $\tau_x$, $x\in X$, i.e.
\[
\tile=\bigcup_{x\in X}\tau_x(\tile).
\]
Hence Proposition~\ref{prop_measure lims propert} item~3 implies that $\mu|_{\tile}$ is the
self-similar measure on $\tile$ by the standard definition of Hutchinson
(\ref{eqn_self_similar_measure}). The measure $\mu$ is the $G$-invariant extension of the
self-similar measure $\mu|_{\tile}$ to the limit $G$-space $\limGs$.


%Let us explain how the measure $\muG$ on $\limGs$ is related to self-similar measures. The tile
%$\tile$ of the limit $G$-space $\limGs$ can be considered as the attractor of the iterated function
%system $\tau_x$, $x\in X$, i.e.
%\[
%\tile=\bigcup_{x\in X}\tau_x(\tile).
%\]
%The measure $\muG$ restricted to $\tile$ satisfies
%\[
%\muG(A)=\sum_{x\in X} \frac{1}{|X|}\muG(\tau_x^{-1}(A))
%\]\notes{maybe normalized?}
%for any Borel set $A\subset\tile$. Hence $\muG|_{\tile}$ is the self-similar measure on $\tile$ by
%the standard definition of Hutchinson (\ref{eqn_self_similar_measure}). The measure $\muG$ is the
%$G$-invariant extension of the self-similar measure $\muG|_{\tile}$ to the limit $G$-space
%$\limGs$.



Let us show how to find the measure of the tile $\tile$. Let $\nucl$ be the nucleus of the action
$(G,X^*)$ identified with its Moore diagram. Replacing each label $(x,y)$ by label $x$ in the
nucleus $\nucl$ we get a strictly right-resolving graph denoted $\gr_{\nucl}$ labeled by elements
of $X$, so that we can apply the methods developed in Section~2.


\begin{theor}\label{theor_tile_measure}
The measure $\mu(\tile)$ is equal to the measure number $\mgr(\gr_\nucl)$ of the nucleus, in
particular it is always integer. Moreover, $\mu(\tile)=1$ if and only if the action satisfies the
open set condition.
\end{theor}
\begin{proof}
By Proposition~\ref{prop_asympt_Moore_diagr} we have
\[
\pi^{-1}_{\limGs[]}(\tile)=\bigcup_{g\in\nucl}F_g\cdot g^{-1},
\]
where the sets $F_g$ are defined using the graph $\gr_\nucl$ (see Section~2). Thus
\[
\mu(\tile)=\sum_{g\in\nucl}\mu_u(F_g)=\mgr(\nucl),
\]
which is integer by Theorem~\ref{theor_measure_graph_integer}. Observe that $\mu(\tile)\geq 1$
because $F_g=\xmo$ for $g=1\in\nucl$.

If the action satisfies the open set condition then $\mu(\tile)=\mgr(\gr_\nucl)=1$ by
Corollary~\ref{cor_meas_sum_meas}.

Suppose now that the action does not satisfy the open set condition. Then there exists an element
$h$ of the nucleus, whose all restrictions are non-trivial. Let $\nucl_1$ be the set of all
restrictions of $h$, and $\gr_{\nucl_1}$ be the corresponding graph. Then by Proposition~\ref{prop
cup B_v eq cup F_v}
\[
\sum_{g\in\nucl_1}\mu_u(F_g)=\mgr(\gr_{\nucl_1})\geq 1.
\]
Thus $\mu(\tile)=\sum_{g\in\nucl}\mu_u(F_g)\geq \mgr(\gr_{\nucl_1})+\mu_u(F_{g=1})\geq 2$.
\end{proof}

\begin{remark}
The measure $\mu(\tile)=\mgr(\gr_{\nucl})$ can be found algorithmically using the remarks after
Proposition~\ref{prop_eigenvect_transition_matr}.
\end{remark}


%\begin{defi}\label{ooo}
%The graph $\nucl_o$ is obtained from $\nucl$ by replacing each label $(x,y)$ by $x$
%\end{defi}
%\begin{defi}
%In the following we will call $\mgr(\nucl_o)$ the measure number of the nucleus $\nucl$ and will
%denote it simply by $\mgr(\nucl)$.
%\end{defi}


%Consider the sets $\mathcal{O}_G=\mathcal{O}(\gr_\nucl)$ and
%$\mathcal{O}_G^*=\mathcal{O}(\gr_\nucl)^*$ defined in Lemma~\ref{lemma_set V open dense}, in other
%words $\mathcal{O}(\gr_\nucl)$ is the set of all $w\in\xmo$ that end in precisely $\muG(\tile)$
%elements of $\nucl$.
%
%\begin{lemma}\label{}
%The sets $\mathcal{O}_G$ is open and dense in $\xmo$, and $\mu_p(\mathcal{O}_G)=1$. For each $w\in
%\mathcal{O}_G$ there is a beginning of $w$ that belongs to $\mathcal{O}_G^*$, and
%$\mathcal{O}_G=\cup_{w\in\mathcal{O}_G^*}\xmo w$.
%
%The sets $\mathcal{O}_G$ and $\mathcal{O}_G\times G$ are closed under the asymptotic equivalence
%relation on $\xmo$ and $\xmo\times G$ respectively.
%\end{lemma}




The next proposition shows that the covering $\limGs=\cup_{g\in G}\tile\cdot g$ is a perfect
multiple covering of multiplicity $\mu(\tile)$.

\begin{prop}\label{points_above}
Every point $x\in\limGs$ is covered by at least $\mu(\tile)$ tiles. The set $\dot{\limGs}$ of all
points $x\in\limGs$, which are covered by exactly $\mu(\tile)$ tiles, is open and dense in
$\limGs$, and its complement has measure $0$.
\end{prop}
\begin{proof}
For each $x\in\limGs$ we define the number $n_x$ of such $g\in G$ that the tile $\tile\cdot g$
contains $x$.

First we prove the inequality. Let $x\in\limGs$ be represented by the pair $w\cdot g$ from
$\xmo\times G$. Then $|V(w)|\geq\mgr(\gr_\nucl)$ by Theorem~\ref{theor_measure_graph_integer}. If
$h\in V(w)$ it means that the sequence $w$ ends in $h$, which by
Proposition~\ref{prop_asympt_Moore_diagr} means that there is a sequence $u_h\in\xmo$ such that $w$
is asymptotically equivalent to $u_h\cdot h$. It follows that $w\cdot g$ is asymptotically
equivalent to $u_h\cdot hg$. It means, in turn, that $x$ belongs to the tile $\tile\cdot hg$ for
every $h\in V(w)$. It follows that $n_x\geq\mgr(\gr_\nucl)$.

Consider the set $\mathcal{O}=\mathcal{O}(\gr_\nucl)$ defined in Lemma~\ref{lemma_set V open dense}
using the graph $\gr_\nucl$, in other words $\mathcal{O}$ is the set of all $w\in\xmo$ that end in
precisely $\mgr(\gr_\nucl)$ elements of $\nucl$. By the same considerations as above we see that if
$w\cdot g$ represents a point $x$ with $n_x=\mgr(\gr_\nucl)$ then $|V(w)|=\mgr(\gr_\nucl)$, that
is, $w$ belongs to $\mathcal{O}$. In other words, the set $\mathcal{O}\times G$ is closed under the
asymptotic equivalence relation on $\xmo\times G$, and it is the inverse image of the set
$\dot{\limGs}$ under the factor map $\pi_{\limGs[]}$. Since the set $\mathcal{O}\times G$ is open
and dense in $\xmo\times G$ by Lemma~\ref{lemma_set V open dense}, the same hold for $\dot{\limGs}$
in $\limGs$. The complement of $\dot{\limGs}$ has measure $0$, because the complement of the set
$\mathcal{O}$ has measure $0$ by Lemma~\ref{lemma_set V open dense}.
\end{proof}

\begin{cor}
Let $\partial\tile$ be the boundary of $\tile$. Then $\mu(\partial\tile)=0$.
\end{cor}
\begin{proof}
Take $x\in\partial\tile$. Suppose that $x$ is covered by exactly $k=\mu(\tile)$ tiles. Let
$U\subset\dot{\limGs}$ be the neighborhood of $x$ which intersects with finitely many tiles, say
with $l\geq k$ tiles. By subtracting any tile that does not contain $x$ we may suppose that $l=k$,
that is, $U$ has nonempty intersection only with those tiles that contain $x$. Now, since $x$ is in
the boundary of $\tile$, there is a point $x'$ in $U$ that is not covered by $\tile$ (otherwise
$x\in U\subset\tile$), thus $x'$ is covered by less than $k$ tiles, contradiction. Hence
$\partial\tile$ is in the complement of $\dot{\limGs}$ and has measure $0$.
\end{proof}

%The preimage of the complement of $\dot{\limGs}$ is the product $\mathcal{O}^c\times G$. Since the complement of
%$\mathcal{O}$, $\mathcal{O}^c$, has measure $0$ by the theorem~\ref{theor_measure_graph_integer}, and $G$ is
%countable, we have that the complement of $\limGs$ has measure $0$.



\begin{remark}\label{rem_meas_inters_tiles}
Two tiles $\tile\cdot g_1$ and $\tile\cdot g_2$ for $g_1,g_2\in G$ have nonempty intersection if
and only if $g_1g_2^{-1}\in\nucl$. Let us show how to find the measure of this intersection. By
Proposition~\ref{prop_asympt_Moore_diagr} we have
\begin{eqnarray*}
\pi^{-1}_{\limGs[]}((\tile\cdot g_1)\cap (\tile\cdot g_2))=\bigcup_{{\scriptsize\begin{array}{c}
                              h_1,h_2\in\nucl \\
                              g_1g_2^{-1}=h_1h_2^{-1}
                            \end{array}}}
(F_{h_1}\cap F_{h_2})\cdot h_1^{-1}g_1,
\end{eqnarray*}
where the sets $F_g$ are defined using the nucleus $\nucl$. The word problem in contracting
self-similar groups is solvable in polynomial time \cite[Proposition~2.13.10]{self_sim_groups} (one
can use the program package \cite{AutomGrp}). The measures of intersections $F_{h_1}\cap F_{h_2}$
can be found for example using method described after Corollary~\ref{cor_meas_gr_rational}.
\end{remark}

%&=&\left(\bigcup_{h_1\in\nucl}F_{h_1}\cdot
%h_{1}^{-1}g_1\right)\bigcap\left(\bigcup_{h_2\in\nucl}F_{h_2}\cdot
%{h_2}^{-1}g_2\right)=\\
%&


The measure $\mu=\mu_{(G,X^{*})}$ on the limit $G$-space $\limGs$ was defined using the specific
self-similar action $(G,X^{*})$ of the group $G$. Let us show that actually this measure depends
only on the associated virtual endomorphism $\phi$, as the limit $G$-space itself. It allows us to
consider the measure space $(\limGs(\phi),\mu)$ independently of the self-similar action. At the
same time, the measure $\mu(\tile)$ may vary for different self-similar actions as the nucleus
does. It is an interesting open question in what cases we can always choose a self-similar action
which satisfies the open set condition (see \cite{mult_an_self_sim,lag_wang:haar} for the abelian
case and applications to wavelets).

\begin{theor}\label{theor_measure_limGs_invariant}
Let $\phi: H\rightarrow G$ be a virtual endomorphism of the group $G$. Let $(G,X^*)$ and $(G,Y^*)$
be the contracting self-similar actions defined respectively by the triples $(\phi,T,C)$ and
$(\phi,T',C')$. Then the homeomorphism $\alpha : \limGs[(G,X^*)]\longrightarrow \limGs[(G,Y^*)]$
from Theorem~\ref{theor_limsX hom limsY} preserves measure, i.e.
\[
\mu_{(G,Y^*)}(\alpha(A))=\mu_{(G,X^*)}(A)
\]
for any Borel set $A$.
\end{theor}
\begin{proof}
Let $\nucl$ be the nucleus and $\tile$ be the tile of the action $(G,X^{*})$. By
Theorem~\ref{theor_tile_measure} and Theorem~\ref{theor_measure_graph_integer}
\begin{equation}\label{eqn_theor_meraXY_mu tile }
\mu_{(G,X^*)}(\tile)=\mgr(\gr_\nucl)=\min_{w\in\xmo} |V(w)|=k\in\mathbb{N},
\end{equation}
where $V(w)$ is defined using the graph $\gr_\nucl$. Consider
\[
\pi_Y^{-1}(\alpha(\tile))=\bigcup_{g\in G} T_g\cdot g\quad (\mbox{here } T_g\subseteq Y^{-\omega}),
\]
where $\pi_Y:Y^{-\omega}\times G\rightarrow \limGs$ is the canonical projection.

Take $w\in Y^{-\omega}$ and let us prove that there exist at least $k$ elements $g\in G$ such that
$w\in T_g$. The tiles $\tile\cdot g$ cover the limit $G$-space $\limGs[(G,X^*)]$, and we can find
$g\in G$ such that if $x=\alpha^{-1}(y)$ for $y=\pi_{Y}(w\cdot g)\in\limGs[(G,Y^*)]$ then $x$
belongs to the tile $\tile$. Then $x$ is represented by the sequence $u\cdot 1$ for $u\in\xmo$.
Equation~(\ref{eqn_theor_meraXY_mu tile }) implies $|V(u)|\geq k$ and so there exist $k$ elements
$h_1,\ldots, h_k\in\nucl$ and sequences $u_1,\ldots,u_k\in\xmo$ such that for every $i$ there
exists a left-infinite path in the nucleus $\nucl$ which ends in $h_i$ and is labeled by $(u,u_i)$.
By Proposition~\ref{prop_asympt_Moore_diagr}
\[
x=\pi_X(u\cdot 1)=\pi_X(u_i\cdot h_i)\quad\mbox{ and }\quad x\cdot h_i^{-1}\in\tile
\]
for every $i=1,\ldots, k$. Then
\[
\pi_Y(w\cdot gh_i^{-1})=\pi_Y(w\cdot g)h_i^{-1}=y\cdot h_i^{-1}=\alpha(x\cdot
h_i^{-1})\in\alpha(\tile)
\]
and thus $w\in T_{gh_i^{-1}}$ for all $i=1,\ldots, k$.

Let $\chi_{T_g}$ be the characteristic function of the set $T_g$. Then $\sum_{g\in
G}\chi_{T_g}(x)\geq k$ for almost all $x$. Integrating we get
$\mu_{(G,Y^*)}(\alpha(\tile))=\sum_{g\in G} \mu_u(T_g)\geq k=\mu_{(G,X^*)}(\tile)$.

Let us now show that $\mu_{(G,Y^*)}(\alpha(\tile_v))\leq\mu_{(G,X^*)}(\tile_v)$ for any $v\in X^*$.
Indeed, let $n=|v|$, then it follows from Theorem \ref{theor_limsX hom limsY} that
$\alpha(\tile_v)=\alpha(\tau_v(\tile))=\tau_u(\alpha(\tile))g$ for some word $u\in Y^n$ and $g\in
G$. Thus
\begin{eqnarray*}
\mu_{(G,Y^*)}(\alpha(\tile_v))&=&\mu_{(G,Y^*)}(\tau_u(\alpha(\tile)))\leq\frac{1}{|X|^n}\mu_{(G,Y^*)}(\alpha(\tile))=\\
&=&\frac{1}{|X|^n}\mu_{(G,X^*)}(\tile)=\mu_{(G,X^*)}(\tile_v).
\end{eqnarray*}

Let us prove that $\mu_{(G,Y^*)}\circ\alpha$ is absolutely continuous with respect to
$\mu_{(G,X^*)}$. Indeed, let $\mu_{(G,X^*)}(A)<\varepsilon$, $\pi^{-1}_{X}(A)=\cup_g T_g\cdot g$.
Then $\sum_g\mu_u(T_g)<\varepsilon$. It follows that there exist $v_{i,g}\in X^*$ such that
$T_g\subset\cup_i\xmo v_{i,g}$ and $\sum_{i,g}|X|^{-|v_{i,g}|}<\varepsilon$. Then
$A\subset\cup_{i,g}\tile_{v_{i,g}}\cdot g$, and we have that
$\sum_{i,g}\mu_{(G,X^*)}(\tile_{v_{i,g}})<\varepsilon\mu_{(G,X^*)}(\tile)$. Then
$\mu_{(G,Y^*)}(\alpha(A))\leq\sum\mu_{(G,Y^*)}(\alpha(\tile_{v_{i,g}}))\leq\sum\mu_{(G,X^*)}(\tile_{v_{i,g}})<\varepsilon\mu_{(G,X^*)}(\tile)$.
Since $\varepsilon$ is arbitrary, we are done.

We will now prove that $\mu_{(G,Y^*)}\circ\alpha\leq\mu_{(G,X^*)}$. Since both
$\mu_{(G,Y^*)}\circ\alpha$ and $\mu_{(G,X^*)}$ are invariant under multiplication by $g\in G$ it
suffices to prove this inequality for sets $A\subset\tile$. Since any Borel $A$ is a union of a
closed set and a set of arbitrarily small measure, it suffices to prove the inequality for closed
sets, as $\mu_{(G,Y^*)}\circ\alpha$ is absolutely continuous with respect to $\mu_{(G,X^*)}$ by
above.

So let $A\subset\tile$ be a closed set. For each $n$, let $A_n$ be the union of all tiles
$\tile_v,v\in X^n,$ that have non-empty intersection with $A$. Let us show that $A=\cap_n A_n$.
Suppose $x\in\cap A_n$, and $x$ is not in $A$. Then for each $n$ there is $v_n\in X^n$ such that
$\tile_{v_n}$ has nonempty intersection with $A$ and $x\in\tile_{v_n}$. It follows that $x$ has
some representation $u_nv_n\in\xmo$. Since the number of such representations is finite, we may
choose subsequence $n_k$ such that $v_{n_k}$ is the beginning of some word $v\in\xmo$. Since $A$ is
compact it follows that $A$ has nonempty intersection with $\cap_k\tile_{v_{n_k}}$, thus
$\cap_k\tile_{v_{n_k}}$ contains at least two points, which is impossible.

Since $A_n$ is the union of some tiles of level $n$, which are disjoint up to sets of measure $0$,
the inequality also holds for all $A_n$. Going to the limit, we get the inequality for $A$.

By interchanging $X$ and $Y$ we get the reverse inequality, and we are done.
\end{proof}

\begin{remark}
It is important in the theorem that we take the uniform Bernoulli measure on $\xmo$. The problem is
that the homeomorphism $\alpha$ may change the Bernoulli measure with a non-uniform weight to a
measure that is not Bernoulli.
\end{remark}


\vspace{0.3cm}\noindent\textbf{Self-similar measure on the limit space $\lims$}. The push-forward
of the uniform Bernoulli measure $\mu_u$ under the factor map $\pi_{\lims[]}:\xmo\rightarrow\lims$
defines the \textit{self-similar measure $\muls$ on the limit space} $\lims$. The measure $\muls$
is a regular Borel probability measure on $\lims$. The shift $\si$ is a measure-preserving
transformation of $\lims$.

Consider the set $\mathcal{U}=\mathcal{U}(\nucl)$ of all sequences $w\in\xmo$ with the property
that every left-infinite path in the nucleus $\nucl$ labeled by $(w,w)$ ends in $1$. Define
$\mathcal{U}^*$ as the set of finite words that satisfy the same condition.



\begin{lemma}\label{lemma U_open_dense_meas1}
The set $\mathcal{U}$ is open and dense in $\xmo$, and $\mu_p(\mathcal{U})=1$. For each $w\in
\mathcal{U}$ there is a beginning of $w$ that belongs to $\mathcal{U}^*$, and
$\mathcal{U}=\cup_{w\in\mathcal{U}^*}\xmo w$.

The sets $\mathcal{U}$ and $\mathcal{U}\times G$ are closed under the asymptotic equivalence
relation on $\xmo$ and $\xmo\times G$ respectively.
\end{lemma}
\begin{proof}
Construct the graph $\gr$ with the elements of $\nucl$ as vertices, for each edge labeled by
$(x,x)$ in $\nucl$ we have the edge in $\gr$ with the same starting and end vertices, and labeled
by $x$; and for each edge in $\nucl$ labeled by $(x,y)$ for $x\neq y$ there is an edge in $\gr$
with the same starting vertex that ends in the trivial element and labeled by $x$. Then the set
$\mathcal{U}$ coincides with the set $\mathcal{O}(\gr)$. Indeed, $h\in V_{\gr}(w)$ if and only if
there is a path in $\nucl$ labeled by $(w,w)$ that ends in $h$, thus $w\in\mathcal{U}$ if and only
if $V_{\gr}(w)=\{1\}$. For every nontrivial element $h\in\nucl$ there exists a word $v\in X^{*}$
such that $h(v)\neq v$. It follows that there exists a path in the graph $\gr$ from $h$ to $1$.
Hence the component $\{1\}$ is the only strongly connected component of the graph $\gr$ without
outgoing edges. By Corollary~\ref{cor_meas_sum_meas} the measure number of $\gr$ is $1$, and
$\mathcal{U}=\mathcal{O}(\gr)$. The first statement of the lemma now follows from
Lemma~\ref{lemma_set V open dense}.


Let us show that the set $\mathcal{U}$ is closed under the asymptotic equivalence relation (then
the set $\mathcal{U}\times G$ is also closed). It is sufficient to show that if there is a path in
$\nucl$ labeled by $(u,v)$ and $u\in\mathcal{U}$ ($v\in\mathcal{U}$) then $v\in\mathcal{U}$
($u\in\mathcal{U}$). Let the path in $\nucl$ labeled by $(u,v)$ end in $h$. It follows that $u$ is
asymptotically equivalent to $v\cdot h$. Suppose there is a path in $\nucl$ labeled by $(v,v)$ that
ends in $g$. It follows that $v$ is asymptotically equivalent to $v\cdot g$. Thus $u$ is
asymptotically equivalent to $v\cdot gh$ which is asymptotically equivalent to $u\cdot h^{-1}gh$.
By definition, there is a path in $\nucl$ labeled by $(u,u)$ which ends in $h^{-1}gh$. Since
$u\in\mathcal{U}$ we get $h^{-1}gh=1$, thus $g=1$.
\end{proof}




%The measure of every tile $\T_v$ for $v\in X^n$ is equal $\muls(\T_v)=\frac{1}{|X|^n}$.

\begin{prop}\label{prop_measure_prop_lims}
Almost every point of $\lims$ has precisely $|X|$ preimages under~$\si$.
\end{prop}
\begin{proof}
Since every point of $\lims$ has at most $|X|$ preimages under $\si$, it is enough to show that for
almost every $w\in\xmo$ the map $\pi_{\lims[]}:\sigma^{-1}(w)\rightarrow\si^{-1}(\pi_{\lims[]}(w))$
is one-to-one. Suppose that for some $w\in\xmo$ and $x\neq y$ in $X$ we have
$\pi_{\lims[]}(wx)=\pi_{\lims[]}(wy)$. It follows that $wx,wy$ are asymptotically equivalent, thus
there is a left-infinite path in $\nucl$ labeled by $(wx,wy)$. It follows that the prefix of this
path labeled by $(w,w)$ must end in the nontrivial element, so $w\not\in\mathcal{U}$. Since
$\mu_u(\mathcal{U})=1$ by Lemma~\ref{lemma U_open_dense_meas1}, we are done.
\end{proof}


%Now we construct a word $u\in X^*$ such that for each vertex $h$ the
%path in $\Gamma'$ labeled by $u$ that start in $h$ ends in the
%trivial element $1$. To do it let $h_1,\ldots,h_s$ be all nontrivial
%elements of $\nucl$, and suppose there is a word $u'$ such that the
%pathes labeled by $u'$ that start from $h_1,\ldots,h_k$ all end in
%$1$. Let $g$ be the end of the path labeled by $u'$ that start in
%$h_{k+1}$. If $g=1$ we are done, since we have the induction step.
%If not, there is $x\in X$ such that $g(x)\neq x$, thus there is an
%edge in $\Gamma'$ from $g$ to $1$, labeled by $x$. If we take the
%word $u'x$, then the pathes from $h_1,\ldots,h_{k+1}$ labeled by
%$u'x$ end in $1$. Note that $V(u)=\{1\}$, thus $\min_{w\in X^*}
%|V(w)|=1$.


The definition of the measure $\muls$ uses encoding of $\lims$ by sequences $\xmo$. At the same
time the space $\lims$ is defined as the space of obits $\limGs/G$. Let us show that we can recover
$(\lims,\muls)$ from the measure space $(\limGs,\mu)$.




%$\mathcal{O}(\nucl)=\mathcal{O}=\mathcal{V}(\nucl_o)$, in other words it is the set of those
%$w\in\xmo$ for which the number of elements of $\nucl$ that are the ends of pathes labeled by $w$
%is minimal. $\mathcal{U}(\nucl)=\mathcal{U}$ is the set of those $w\in\xmo$ for which whenever
%there is a path in $\nucl$ labeled by $(w,w)$, it must end in $1$. $\mathcal{O}^*$
%($\mathcal{U}^*$) is defined in the same way as $\mathcal{O}$ ($\mathcal{U}$), with $X^*$ instead
%of $\xmo$.


%\begin{lemma}\label{o_u}
%The sets $\mathcal{O}$, $\mathcal{U}$ are open and dense in $\xmo$, and have measure $\mu_p$ equal
%to $1$. For each $w\in\mathcal{O}$ ($\mathcal{U}$) there is a beginning of $w$ that belongs to
%$\mathcal{O}^*$ ($\mathcal{U}^*$). Equivalently, $\mathcal{O}=\cup_{w\in\mathcal{O}^*}\xmo w$
%($\mathcal{U}=\cup_{w\in\mathcal{U}^*}\xmo w$).
%\end{lemma}
%\begin{proof}
%For the set $\mathcal{O}$ the statements of the lemma follow from the lemma~\ref{lemma_o}, since
%$\mathcal{O}(\nucl)=\mathcal{V}(\nucl_o)$.
%
%For the set $\mathcal{U}$ consider the graph $\nucl_u$ with the elements of $\nucl$ as vertices,
%for each edge labeled by $(x,x)$ in $\nucl$ we have the edge in $\nucl_u$ with the same starting
%and end vertices, and labeled by $x$; and for each edge in $\nucl$ labeled by $(x,y)$, $x\neq y$
%there is an edge in $\nucl_u$ with the same starting vertex that ends in the trivial element and
%labeled by $x$. Then $\mathcal{U}(\nucl)=\mathcal{V}(\nucl_u)$. Indeed, $h\in V_{\nucl_u}(w)$ iff
%there is a path in $\nucl$ labeled by $(w,w)$ that ends in $h$, thus $w\in\mathcal{U}$ if and only
%if $V_{\nucl_u}(w)=\{1\}$. Since $\nucl_u$ satisfies the open set condition, it follows that
%$\mathcal{U}(\nucl)=\mathcal{V}(\nucl_u)$ by corrollary~\ref{ops}, and the statements of the lemma
%for $\mathcal{U}$ follow from the Lemma~\ref{lemma_o}.
%\end{proof}
%
%\begin{lemma}\label{o_u_asymp_x}
%The sets $\mathcal{O}$ and $\mathcal{U}$ are closed under the asymptotic equivalence relation on
%$\xmo$. The sets $\mathcal{O}\times G$ and $\mathcal{U}\times G$ are closed under the asymptotic
%equivalence relation on $\xmo\times G$.
%\end{lemma}
%\begin{proof}
%It is sufficient to show that if there is a path in $\nucl$ labeled by $(u,v)$ and
%$u\in\mathcal{U}$ ($v\in\mathcal{U}$) then $v\in\mathcal{U}$ ($u\in\mathcal{U}$). Let the path in
%$\nucl$ labeled by $(u,v)$ end in $h$. It follows that $u$ is asymptotically equivalent to $v\cdot
%h$. Suppose there is a path in $\nucl$ labeled by $(v,v)$ that ends in $g$. It follows that $v$ is
%asymptotically equivalent to $v\cdot g$. Thus $u$ is asymptotically equivalent to $v\cdot gh$ which
%is asymptotically equivalent to $u\cdot h^{-1}gh$. By definition, there is a path in $\nucl$
%labeled by $(u,u)$ which ends in $h^{-1}gh$. Since $u\in\mathcal{U}$ we get $h^{-1}gh=1$, thus
%$g=1$.
%
%\end{proof}



\begin{prop}\label{x_j_homeo}
Let $\rho$ be the factor map $\limGs\rightarrow\lims$. Then for any $u\in\mathcal{U}^*$ the
restriction $\rho|_{\tile_u}:\tile_u\rightarrow\T_u$ is a homeomorphism,
$\rho^{-1}(\T_u)=\sqcup_{g\in G}\tile_u\cdot g$.
\end{prop}
\begin{proof}
By definition, the map $\rho|_{\tile_u}:\tile_u\rightarrow\T_u$ is surjective, and $\tile_u$ is
compact. Hence in order to prove that $\rho$ is a homeomorphism it is left to show that $\rho$ is
injective on $\tile_u$. Take $x,y\in\tile_u$, and let $wu\in\xmo$ represent $x$ and $vu\in\xmo$
represent $y$. Suppose that $\rho(x)=\rho(y)$. By Proposition~\ref{prop_asympt_Moore_diagr} it
means that there is a left-infinite path in $\nucl$ labeled by $(wu,vu)$. Since
$u\in\mathcal{U}^*$, this path must end in $1$. It follows, that $wu$ and $vu$ represent the same
point of the tile $\tile_u$, and $x=y$.

To prove the second claim, take $x\in\tile_u\cdot g\cap\tile_u\cdot g'$. Then $x$ is represented by
two asymptotically equivalent sequences $wu\cdot g$ and $w'u\cdot g'$. It follows that there is a
path in the nucleus $\nucl$ labeled by $(u,u)$ which ends in $g'g^{-1}$. Then $g=g'$ since
$u\in\mathcal{U}^*$.
\end{proof}



\begin{theor}
The projection $\limGs\rightarrow\lims$ is a covering map up to sets of measure zero.
\end{theor}
\begin{proof}
Consider the sets $\tilde{\limGs}=\pi_{\limGs[]}(\mathcal{U}\times G)$ and
$\tilde{\lims}=\pi_{\lims[]}(\mathcal{U})$. It follows that $\tilde{\limGs}/G=\tilde{\lims}$. Since
the set $\mathcal{U}$ has measure $1$ by Lemma~\ref{lemma U_open_dense_meas1}, the complements of
$\mathcal{U}\times G$, of $\tilde{\limGs}$, and of $\tilde{\lims}$ have measure $0$.

%Denote by $\tilde{\limGs}$ the image of $\mathcal{U}\times G$ under the canonical projection
%$\xmo\times G\rightarrow\limGs$, and by $\tilde{\lims}$ the image of $\mathcal{U}$ under the
%canonical projection $\xmo\rightarrow\lims$.

Since the group $G$ acts properly on $\limGs$, the same holds for $\tilde{\limGs}$. It is left to
prove the freeness. Suppose $x\cdot g=x$ for $x\in\tilde{\limGs}$ and $g\in G$. Let $u\cdot h$ be a
representative of $x$ in $\xmo\times G$. It follows that $u\cdot hg$ is asymptotically equivalent
to $u\cdot h$, thus there is a path in $\nucl$ labeled by $(u,u)$ that ends in $h^{-1}gh$. Since
$u\in\mathcal{U}$ we have $g=1$.

Hence the projection $\tilde{\limGs}\rightarrow\tilde{\lims}$ is a covering map, and the statement
follows.
\end{proof}


Let $(\limGs[],\nu)$ be a locally compact measure space, the group $G$ acts freely and properly
discontinuously on $\limGs[]$ by homeomorphisms, and the measure $\nu$ is $G$-invariant. There is a
unique measure $\nu_{*}$ on the quotient space $\limGs[]/G$, called the \textit{quotient measure},
with the property that if $U$ is an open subset of $\limGs[]$ such that $Ug\cap Uh=\emptyset$ for
all $g,h\in G$, $g\neq h$, then $\nu_{*}(U/G)=\nu(U)$.

%The projection of the measure $\mu$ on $\lims$ is equal to the measure $\muls$.

\begin{prop}
The quotient measure $\mu_{*}$ of the limit $G$-space $(\limGs,\mu)$ coincides with the measure
$\muls$ on the limit space $\lims$.
\end{prop}
\begin{proof}
Consider the sets $\tilde{\lims}$ and $\tilde{\limGs}$ of full measure from the previous theorem.
For every $u\in\mathcal{U}^*$ we have $\T_u\subset\tilde{\lims}$, $\tile_u\cdot
g\subset\tilde{\limGs}$ for every $g\in G$, and $\rho^{-1}(\T_u)=\bigsqcup_{g\in G} \tile_u\cdot g$
by Proposition~\ref{x_j_homeo}. Since $\tilde{\lims}=\cup_{u\in\mathcal{U}^*}\T_u$ by
Lemma~\ref{lemma U_open_dense_meas1} it suffices to show that for any $u\in\mathcal{U}^*$ the
restriction of $\mu_{*}$ on $\T_u$ is equal to the restriction of $\muls$. Take a Borel set
$A\subset\tile_u$ and consider its preimage $\pi^{-1}_{\limGs[]}(A)=\sqcup_{g\in\nucl} A_g\cdot g$.
Then $\pi^{-1}_{\lims[]}(\rho(A))=\sqcup_{g\in\nucl} A_g$. Here the union is disjoint because if
$w\in{A_g\cap{}A_h}$ then $w\cdot g$ and $w\cdot h$ are asymptotically equivalent to $vu$ and $v'u$
respectively, for some $v,v'\in\xmo$. It follows that $vu\cdot g^{-1}$ is equivalent to $v'u\cdot
h^{-1}$. Since $u\in\mathcal{U}$, $g=h$. Hence
\[
\muls(\rho(A))=\mu_u(\pi^{-1}_{\lims[]}(\rho(A)))=\mu_u(\sqcup_{g\in\nucl}
A_g)=\sum_{g\in\nucl}\mu_u(A_g).
\]
By the property of the quotient measure
\[
\mu_{*}(\rho(A))=\mu(A)=\sum_{g\in\nucl}\mu_u(A_g)=\muls(\rho(A)).
\]
\end{proof}

%\begin{lemma}
%Let $A\subset\limGs$. If $\pi^{-1}(A)=\cup_g A_g\cdot g$ then $\pi^{-1}(\rho(A))=\cup_g A_g$. If
%$A\subset\tilde{\limGs}$ then the last union is disjoint.
%\end{lemma}
%\begin{proof}
%Let $w\in A_g$. Then $\pi(wg)\in A$, thus $\rho(\pi(wg))=\pi(w)\in\rho(A)$, thus $w\in\pi^{-1}(\rho(A))$.
%Conversely, let $w\in\pi^{-1}(\rho(A))$. Then $\pi(w)=\rho(\pi(vg))$ for some $v\in A_g$. Thus $\pi(w)=\pi(v)$,
%so $v$ is assymptotically equivalent to $wh$ for some $h\in\nucl$. It follows that $vg$ is assymptotically
%equivalent to $whg$, so $\pi(whg)=\pi(vg)\in A$, thus $w\in A_{hg}$.
%  Suppose $A\subset\tile\otimes u$, and $w\in A_g\cap A_{g'}$. Then $x=\pi(wg)=\pi(wg')\in\tilde{\limGs}$. We have that $x\cdot g^{-1}g'=x$. By the lemma~\ref{raz2} we have that $g=g'$.
%\end{proof}
%
% We just need to prove that $\mu_\si(\rho(A))=proj_*\mu(\rho(A))$. By definition
%of $\mu_\si$, $\mu_\si(\rho(A))=\mu(\pi^{-1}(\rho(A)))=\mu(\cup_g A_g)=\sum_g\mu(A_g)$. By the main property of
%$proj_*\mu$, $proj_*\mu(\rho(A))=\mu(A)=\mu(\pi^{-1}(A))=\sum_g\mu(A_g)$.

%Let $u\in X^*$ be a finite word such that whenever there is a path in $\nucl$ labeled by $(u,u)$ it
%must end in $1$. Then clearly $\xmo u\subset\mathcal{U}$, $\T_u\subset\tilde{\lims}$ and
%$\tile_u\cdot g\subset\tilde{\limGs}$. By the theorem~\ref{x_j_homeo} we have that
%\[
%\rho^{-1}(\T_u)=\bigsqcup_{g\in G} \tile_u\cdot g,
%\]
%where the union is disjoint.

%Note that $\rho$ is a homeomorphism from $\tile_u$ to $\T_u$, because the action of $G$
%on $\tilde{\limGs}$ is free.  Indeed, suppose $x,y\in\tile_u$, $x=yg$ for some $g\in G$. Say $x=\pi(wu)$,
%$y=\pi(w'u)$. It follows that $wu$
%is assymptotically equivalent to $w'ug$, and by the definition of $u$ $g=1$.

%Indeed, if $x\in\tile_u\cdot g$ then $x=\pi(wug)$ for some $w\in\xmo$. Thus $\rho(x)=\pi(wu)\in\T_u$. Conversly,
%let $x\in\tilde{\limGs}$ such that $\rho(x)\in\T_u$. Let $x=\pi(w'g)$, $\rho(x)=\pi(wu)$. Then also
%$\rho(x)=\pi(w')$ thus $w'$ is assymptotically equivalent to $wuh$ for some $h\in\nucl$. Thus $w'g$ is
%assymptotically equivalent to $wu hg$, and so $x=\pi(w'g)=\pi(wu hg)\in\tile\otimes uhg$.


\begin{cor}
Theorem~\ref{theor_measure_limGs_invariant} holds for the limit space $\lims$, i.e. the measure
space $(\lims,\muls)$ depends only on the associated virtual endomorphism.
\end{cor}

\begin{cor}
$\muls(\T_u)=\mu(\tile_u)=\frac{1}{|X|^{|u|}}\mgr(\tile)$ for $u\in\mathcal{U}^*$.
\end{cor}




%
%
%Let us define the branches $\si_x$ for $x\in X$ of the inverse of the shift $\si$. Notice that
%every point of the limit space $\lims$ is contained in $\leq |\nucl|$ tiles of the $n$-th level for
%every $n$. Take a point $t\in\lims$. Let $A_n$ be the union of all tiles $\T_u$ of the $n$-th level
%that contain the point $t$ (here $t=\cap_n A_n$). Let $B_n$ be the union of the corresponding tiles
%$\T_{ux}$. The sequence $B_n$ is a decreasing sequence of nonempty compact sets, hence the
%intersection $\cap_n B_n$ is nonempty and contains one point, because the diameters of tiles goes
%to zero. We define $\si_x(t)=\cap_n B_n$. Clearly, $\si(\si_x(t))=t$.
%
%\red{ Consider\notes{remove this?} the branches $\si_x$ for $x\in X$ of the inverse of $\si$. Then
%\begin{equation}\label{eqn_self_similar_measure_muls}
%\muls(A)=\sum_{x\in X}\frac{1}{|X|}\muls(\si_x^{-1}(A))
%\end{equation}
%for any Borel set $A$, which follows from analogous formula for the space $(\xmo,\sigma,\mu_u)$. }


\begin{theor}\label{theor_lims_conju_Bernoulli}
$(\lims,\si,\muls)$ is conjugate to the one-sided Bernoulli $|X|$-shift.
\end{theor}
\begin{proof}
We will use notions and results from \cite{hoffman:Bernoulli}. First recall that a measure
preserving map with entropy $\log{d}$ is called {\it uniformly $d$-to-one endomorphism} if it is
almost everywhere $d$-to-one and the conditional expectation of each preimage is $1/d$.  The
standard example is $(\xmo,\sigma,\mu_u)$, which is uniformly $|X|$-to-one. Next we want to show
that $(\lims,\si,\muls)$ is also uniformly $|X|$-to-one. By
Proposition~\ref{prop_measure_prop_lims} the map $\pi:\xmo\rightarrow\lims$ is injective on the
preimages $\sigma^{-1}(w)$ for almost all $w\in\xmo$, that is, $\pi$ is {\it tree adapted} in
terminology of \cite{hoffman:Bernoulli}. We can apply Lemma~2.3 from \cite{hoffman:Bernoulli},
which says that a tree adapted factor of a uniform $d$-to-one endomorphism is again uniform
$d$-to-one endomorphism. In particular, the shift $\si$ is the map of maximal entropy $\log |X|$.

To prove the theorem, we use the following Theorem~5.5 in \cite{hoffman:Bernoulli}.

\begin{theor}
A uniform $d$-to-one endomorphism $(Y,S,\mu)$ is one-sidedly conjugated to the one-sided Bernoulli
$d$-shift if and only if there exists a generating function $f$ so that $(Y,S,\mu)$ and $f$ are
tree very weak Bernoulli.
\end{theor}

%and show that $(\lims,\si,\mu_\si)$ is tree very weak Bernoulli for some generating $f$.
%, $\mu$ be a probability measure,

Recall the definition of tree very weak Bernoulli and generating function. Let $(Y,S,\mu)$ be
uniformly $d$-to-one and $f:Y\rightarrow{R}$ be a tree adapted function to a compact metric space
$R$ with metric $D$. The function $f$ is called \textit{generating} if the $\sigma$-algebra on $Y$
is generated by $S^{-i}f^{-1}(\mathcal{B}),i\geq{0}$, where $\mathcal{B}$ is the $\sigma$-algebra
of Borel sets of the space $R$.

Informally, ``tree very weak Bernoulli'' means that for almost all pairs of points in $Y$ their
trees of preimages are close. To give a formal definition note that since $S$ is uniformly
$d$-to-one, for almost all points $y\in Y$ the set $S^{-k}(y)$ contains exactly $d^k$ points, i.e.
the tree of preimages is a $d$-regular rooted tree. The set $\{1,\ldots, d\}^{*}$ of finite words
over $\{1,\ldots,d\}$ can be considered as a $d$-regular rooted tree, where every word $v$ is
connected with $vx$ for $x\in\{1,\ldots,d\}$, and the root is the empty word $\emptyset$. We can
use the tree $\{1,\ldots, d\}^{*}$ to label the trees of preimages. For almost all points $y\in Y$
there is a map $T_y:\{1,\dots,d\}^*\rightarrow{Y}$ such that $T_y(\emptyset)=y$ and
$T_y(\sigma(v))=S(T_y(v))$ for all nonempty words $v\in \{1,\dots,d\}^*$. Every map $T_y$ is tree
adapted, and it is uniquely defined up to an automorphism of the tree $\{1,\dots,d\}^*$. Then
$(Y,S,\mu)$ and $f$ are called {\it tree very weak Bernoulli} if for any $\varepsilon>0$ and all
sufficiently large $n$ there is a set $W=W(\varepsilon,n)\subset Y$ with $\mu(W)>1-\varepsilon$
such that for any $y,y'\in{W}$
\begin{equation}
\label{} t_n(y,y')=\min_\psi \frac{1}{n}\sum_{v\in \{1,\dots,d\}^{*}, |v|\leq
n}d^{-|v|}D(f(T_y(v)),f(T_{y'}(\psi(v))))<\varepsilon,
\end{equation}
where the minimum is taken over all automorphisms $\psi$ of the tree $\{1,\dots,d\}^*$. Notice that
the definition of $t_n$ does not depend on the choice of $T_y$.

%\begin{lemma}\label{lemmaodin}
%The canonical projection $\pi:\xmo\rightarrow\lims$ is tree-adapted.
%\end{lemma}
%\begin{proof}
%Suppose that for some $w\in\xmo$ and $x\neq y$ in $X$ we have
%$\pi(wx)=\pi(wy)$. It follows that $wx,wy$ are asymptotically
%equivalent, thus there is a left-infinite path in $\nucl$ labeled by
%$(wx,wy)$. It follows that the subpath labeled by $(w,w)$ must end
%in the nontrivial element.
%
%For $w\in\xmo$ let $V(w)$ be the set of those elements $h$ of
%$\nucl$ that there exists a path in $\nucl$ labeled by $(w,w)$ that
%ends in $h$. Clearly $V(w)$ contains the trivial element. To prove
%the lemma it suffices to show that the set of those $w$ that
%$|V(w)|=1$ has measure $1$. To do that, construct a graph $\Gamma'$
%with the elements of $\nucl$ as vertices, for each edge labeled by
%$(x,x)$ in $\nucl$ we have the edge in $\Gamma'$ with the same
%beginning and end vertices, and labeled by $x$; and for each edge in
%$\nucl$ labeled by $(x,y)$, $x\neq y$ there is an edge in $\Gamma'$
%with the same beginning that ends in the trivial element and labeled
%by $x$. Then $V(w)$ is also the set of those $h$ that there is a
%path in $\Gamma'$ that ends in $h$ labeled by $w$. By construction,
%there exists a path in $\Gamma'$ from any vertex to the trivial
%state (open set condition for graphs). Hence we can construct a word
%$u\in X^{*}$ such that for each vertex $h$ the path in $\Gamma'$
%labeled by $u$ that start in $h$ ends in the trivial element $1$.
%Note that $V(u)=\{1\}$, thus $\min_{w\in X^*} |V(w)|=1$.
%
%Now, by Theorem~\ref{theor_measure_graph_integer} $\min_{w\in\xmo}
%|V(w)|=\min_{w\in X^*} |V(w)|=1$. By the proof of it, the set of
%those $w\in\xmo$ for which the minimum $\min_{w\in\xmo} |V(w)|$ is
%attained, that is, for which $|V(w)|=1$, is of measure $1$.
%%
%%Now we construct a word $u\in X^*$ such that for each vertex $h$ the
%%path in $\Gamma'$ labeled by $u$ that start in $h$ ends in the
%%trivial element $1$. To do it let $h_1,\ldots,h_s$ be all nontrivial
%%elements of $\nucl$, and suppose there is a word $u'$ such that the
%%pathes labeled by $u'$ that start from $h_1,\ldots,h_k$ all end in
%%$1$. Let $g$ be the end of the path labeled by $u'$ that start in
%%$h_{k+1}$. If $g=1$ we are done, since we have the induction step.
%%If not, there is $x\in X$ such that $g(x)\neq x$, thus there is an
%%edge in $\Gamma'$ from $g$ to $1$, labeled by $x$. If we take the
%%word $u'x$, then the pathes from $h_1,\ldots,h_{k+1}$ labeled by
%%$u'x$ end in $1$. Note that $V(u)=\{1\}$, thus $\min_{w\in X^*}
%%|V(w)|=1$.
%\end{proof}

Let us show that $(\lims,\si,\muls)$ is tree very weak Bernoulli for the identity map
$id:\lims\rightarrow\lims$. It immediately follows that $id$ is tree adapted and generating. Take a
point $x\in\lims$ and let $x$ be represented by some $w\in\xmo$. Define the map
$T_x:X^{*}\rightarrow\lims$ by the rule $T_x(v)=\pi(wv)$. It is enough to show that $
t_n(x,x')\rightarrow 0, n\rightarrow\infty, $ for almost all $x,x'\in\lims$. Using
(\ref{eqn_diam_tile_goes_zero}) we can find $n_1$ such that $\max_{v\in
X^n}\rm{diam}(\T_v)<\epsilon/2$ for all $n\geq n_1$. It means that $D(\pi(wv),\pi(w'v))<\epsilon/2$
for all $v,|v|\geq n_1$ (here $D$ is a fixed metric on the limit space $\lims$). Thus, taking
$\psi$ to be the identical tree automorphism, we have that
\[
t_n(x,x')<\frac{n_1}{n} \rm{diam}(\lims)+\frac{n-n_1}{n}\frac{\epsilon}{2}<\frac{n_1}{n}
\rm{diam}(\lims)+\frac{\epsilon}{2}<\epsilon
\]
for $n>2n_1\rm{diam}(\lims)/\epsilon$. Hence $(\lims,\si,\muls)$ and $id$ are tree very weak
Bernoulli, which finishes the proof.
\end{proof}


%\begin{lemma}
%$(\lims,\si,\mu_\si)$ and $id$ are tree very weak Bernoulli.
%\end{lemma}
%\begin{proof}
%%The map $id$ is tree adapted because $(\lims,\si,\mu_\si)$ is
%%uniform $|X|$-to-one endomorphism.
%If $\zeta$ is represented by the sequence $w\in\xmo$ then $T_\zeta=\{\pi(wv)|v\in X^*\}$.
%\end{proof}


%(measure-preserving map with entropy $\log |X|$ which is almost
%everywhere $|X|$-to-one and for which the conditional expectation of
%each preimage is precisely $1/|X|$)



In the same way we introduce the measure $\muls_\e$ on the limit solenoid $\sol$ as the
push-forward of the uniform Bernoulli measure on $\xz$. It is easy to see that $(\sol,\e,\muls_\e)$
is the inverse limit of dynamical systems $(\lims,\si,\muls)$ (see \cite[page~27]{dyn_sysII}). In
particular, we get


\begin{cor}
$(\sol,\e,\muls_\e)$ is conjugate to the two-sided Bernoulli $|X|$-shift.
\end{cor}




\section{Applications and Examples}\label{section_Applictions}

\textbf{Invariant measures on nilpotent Lie groups.} Let $G$ be a finitely generated torsion-free
nilpotent group. Let $\phi:H\rightarrow G$ be a contracting surjective virtual endomorphism such
that the associated self-similar action is faithful (i.e. $\phi$-core($H$) is trivial in the
terminology of \cite{sidki:virt_nilp}). Then $\phi$ is also injective by Theorem~1 in
\cite{sidki:virt_nilp}, thus $\phi$ is an isomorphism and we can apply Theorem~6.1.6 from
\cite{self_sim_groups}. The group $G$ and its subgroup $H$ are uniform lattices of a simply
connected nilpotent Lie group $L$ by Malcev's completion theorem. The isomorphism
$\phi:H\rightarrow G$ extends to a contracting automorphism $\phi_L$ of the Lie group $L$. There
exists a $G$-equivariant homeomorphism $\Phi:\limGs\rightarrow L$ such that
$\phi_L(\Phi(t))=\Phi(\tau_{x_0}(t)\cdot g_0)$ for every $t\in\limGs$ and fixed $x_0\in X$ and
$g_0\in G$.


%The push-forward $\Phi_*\mu$ of the measure $\mu$ on the limit $G$-space $\limGs$ is the (right)
%Haar measure on the Lie group $L$, because it is a non-zero regular Borel measure on $L$ that is
%translation invariant.


% such that $\Phi(\zeta\otimes\phi(g_1)g_2)=\phi_L(\Phi(\zeta)g_1)g_2$.

%Let $G$ be a finitely generated virtually nilpotent
%group. Let $\phi:H\rightarrow G$ be contracting surjective virtual endomorphism such that the associated
%self-similar action $(G,X^{*})$ is faithful (i.e. $\phi$-core($H$) is trivial in the terminology of \cite{}).
%There is a normal nilpotent torsion-free subgroup of finite index $G_1\lhd G$ such that $\phi:H_1=G_1\cap
%H\rightarrow G_1$ is surjective virtual endomorphism of $G_1$ (Lemma~6.1.10 in \cite{}). Since the action is
%faithful, the map $\phi|_{H_1}$ is also injective by Theorem~1 in \cite{} and thus $\phi|_{H_1}$ is an
%isomorphism.


%Recall theorem 6.1.6 from \ref{Nekrashevych}:
%\begin{theor}
%Let $\phi:G\rightarrow G$ be a surjective and injective contracting
%virtual endomorphism and let $(G,X)$ be the associated self-similar
%action. Then there exist a nilpotent connected and simply connected
%Lie group $L$, a co-compact proper action of $G$ on $L$ by affine
%transformations and a $G-$~equivariant homeomorphism
%$\Phi:\limGs\rightarrow L$. Moreover, thevirtual endomorphism $\phi$
%is induced by a contracting automorphism $\phi_L$ of the Lie group
%$L$ and $\Phi(\zeta\otimes\phi(g_1)g_2)=\phi_L(\Phi(\zeta)g_1)g_2$.
%\end{theor}


\begin{prop}\label{haar}
The push-forward $\Phi_*\mu$ of the measure $\mu$ on the limit $G$-space $\limGs$ is the (right)
Haar measure on the Lie group $L$.
%
% Let $\mu$ be the measure on $\limGs$ induced by the uniform measure on $X$. Then push-forward $\Phi_*\mu$ is
%the (right) Haar measure on $L$.
\end{prop}
\begin{proof}
The measure $\Phi_*\mu$ is a non-zero regular Borel measure on $L$. It is left to prove that it is
translation invariant. Since the measure $\mu$ is $G$-invariant and the map $\Phi$ is
$G$-equivariant, the measure $\Phi_*\mu$ is $G$-invariant. By the property of the map $\Phi$ we
have $\phi_L(\Phi(A))=\Phi(\tau_{x_0}(A))g_0$ for every Borel set $A\subset\limGs$. Notice that
since the map $\phi$ is injective, we get $\mu(A)=|X|\mu(\tau_x(A))$ (see the proof of
Proposition~\ref{prop_measure lims propert}) and hence $\Phi_*\mu(B)=|X|\Phi_*\mu(\phi_L(B))$ for
every Borel set $B\subset L$. It follows that the measure $\Phi_*\mu$ is $\cup_n
\phi_L^n(G)$-invariant. Since $\phi_L$ is contracting, the set $\cup_n \phi_L^n(G)$ is dense in the
Lie group $L$. Hence $\Phi_*\mu$ is $L$-invariant.
\end{proof}

%Let $L'$ be the subgroup of elements $l$ of $L$ that leave $\Phi_*\mu$ invariant, that is
%$\Phi_*\mu(Bl)=\Phi_*\mu(B)$ for any $B\subset L$. It is shown in the proof of the theorem 6.1.6
%that there is a subgroup $H$ of $G$ which is also a uniform lattice in $L$, and $\Phi(\zeta
%h)=\Phi(\zeta)h$. Since $\mu$ is invariant with respect to right multiplication of the elements of
%$H$ and $\Phi$ is $G-$~equivariant, we have that $H\subset L'$. Take $g_1=g_2=1$ in theorem 6.1.6,
%then we have that $\Phi(A\otimes\phi(1)1)=\phi_L(\Phi(A))$ for $A\subset\limGs$. Take $m=\phi(1)1$
%in the proposition~\ref{measure_invariant}, we have $\mu(A)=d\mu(A\otimes\phi(1)1)$. It follows
%that $\Phi_*\mu(B)=d\Phi_*\mu(\phi_L(B))$ for any $B\subset L$. Thus $l\in L'$ if and only if
%$\phi_L(l)\in L'$; so $\cup_n \phi_L^n(H)\subset L'$. Since $\phi_L$ is contracting, it follows
%that $L'$ is dense in $L$. To prove that it is equal, note that a regular Borel measure is uniquely
%defined by its action on the compactly supported continuous functions. For any compactly supported
%continuous $f$ on $L$, and any sequence $l_j$ of elements $L'$ that tends to $l$ we have
%$$\int fd\Phi_*\mu=\int fd (l_j)_*\Phi_*\mu=\int f(xl_j)d\Phi_*\mu(x)\rightarrow \int f(xl)d\Phi_*\mu(x)$$
%So $L'$ is closed in $L$, and thus equal to it.


The same observation holds in a more general setting of finitely generated virtually nilpotent
groups (or under the conditions of Theorem~6.1.6).


\vspace{0.3cm}\noindent\textbf{Lebesgue measure of self-affine tiles.} Let $A$ be an $n\times n$
integer expanding matrix, where expanding means that every eigenvalue has modulus $>1$. The lattice
$\Z^n$ is invariant under $A$, and we can choose a coset transversal $D=\{d_1,\ldots,d_m\}$ for
$\Z^n/A(\Z^n)$, where $m=|det(A)|$.
%The set $D$ is called the \textit{digit set}.
There exists a unique nonempty compact set $T=T(A,D)\subset\mathbb{R}^n$, called (standard)
\textit{integral self-affine tile}, satisfying
\[
A(T)=\bigcup_{d\in D} T+d.
\]
The tile $T$ has positive Lebesgue measure, is the closure of its interior, and the union above is
nonoverlapping (the sets have disjoint interiors) \cite{lag_wang:self_affine}. It is well known
\cite{lag_wang:int_self_affine_I} that the tile $T$ has integer Lebesgue measure. The question how
to find this measure is studied in
\cite{lag_wang:int_self_affine_I,he_hui:self_affine,deng_he:int_self_affine}, and finally answered
in \cite{gab_yu:natur_til}. A related question how to find the measure of intersection $T\cap
(T+a)$ for $a\in\mathbb{Z}^n$ is studied in \cite{gab_yu:natur_til,int_two_tiles}. Let us show how
to answer these questions using the theory of self-similar groups.

The inverse of the matrix $A$ can be considered as the contracting virtual endomorphism $A^{-1}:
A(\Z^n)\rightarrow\Z^n$ of the group $\Z^n$, which is actually an isomorphism so that we can apply
the previous example of this section. Put $X=\{x_1,\ldots,x_m\}$ and let $(\mathbb{Z}^n,X^{*})$ be
the self-similar contracting action defined by the virtual endomorphism $A^{-1}$, the coset
transversal $D$, and the trivial cocycle $C=\{1,\ldots,1\}$ (see Section~3). The group $\Z^n$ is
the uniform lattice in the Lie group $\mathbb{R}^n$. Hence by Theorem~6.1.6 in
\cite{self_sim_groups} (see also Section~6.2 there) there exists a $\Z^n$-equivariant homeomorphism
$\Phi:\limGs[\Z^n]\rightarrow\mathbb{R}^n$ given by
\[
\Phi(\ldots x_{i_2} x_{i_1}\cdot g)=g+A^{-1}d_{i_1}+A^{-2}d_{i_2}+\ldots
\]
for $i_j\in\{1,\ldots,m\}$ and $g\in\Z^n$. The image of the tile $\tile$ is the self-affine tile
$T$.


%Let $(\mathbb{Z}^n,X)$ be the faithful self-similar recurrent contracting action of $\mathbb{Z}^n$.Let
%$X=\{x_0,\dots,x_{d-1}\}$. Then by \cite{Nekrashevych} there is an integer expanding (that is, every eigenvalue
%in absolute value is greater then $1$) matrix $A$ of the size $n$, and a coset transversal of $A(\mathbb{Z}^n)$
%in $\mathbb{Z}^n$ $\{r_0,\dots,r_{d-1}\}$, where $r_{x_0}=0$, such that $g\cdot x=y\cdot h$ if and only if
%$g+r_x=r_y+Ah$, for $g,h\in\mathbb{Z}^n$, and $x,y\in X$. Moreover, the map $\Phi:\ldots z_2 z_1 g\mapsto
%g+A^{-1}r_{z_1}+A^{-2}r_{z_2}+\ldots$ from $\limGs$ to $\mathbb{R}^n$ is a $\mathbb{Z}^n$-equivariant
%homeomorphism.
%
%
%Let $\mu$ be the regular Borel measure on $\limGs$ that is induced
%by the uniform Bernoulli measure on $\xmo$, as described in
%section~\ref{Ssmls}. Let $\nu=\Phi_*\mu$ be the push-forward of
%$\mu$ under $\Phi$.

\begin{prop}
The push-forward $\Phi_*\mu$ of the measure $\mu$ on the limit $G$-space $\limGs[\Z^n]$ is the
Lebesgue measure $\lambda$ on $\mathbb{R}^n$.
\end{prop}
\begin{proof}
The measure $\Phi_*\mu$ is the Haar measure on $\mathbb{R}^n$ by the above example. Since the Haar
measure is unique up to multiplicative constant, we have that $\Phi_*\mu=c\lambda$ for some
constant $c>0$ and we need to prove that $c=1$.

Recall that $\mu(\tile)$ is integer by Proposition~\ref{theor_tile_measure}, and almost every point
of $\limGs$ is covered by $\mu(\tile)$ tiles $\tile\cdot g$ by Proposition~\ref{points_above}.
Hence $\Phi_*\mu(T)=\Phi_*\mu(\Phi(\tile))=\mu(\tile)$ is integer and almost every point of
$\mathbb{R}^n$ is covered by $\mu(\tile)$ tiles $T+g$, $g\in\Z^n$, with respect to the measure
$\Phi_*\mu$, and thus with respect to the Lebesgue measure $\lambda$. It follows that, if
$\chi_{T+g}$ is the characteristic function of $T+g$, then $\sum_{g\in\Z^n}\chi_{T+g}=\mu(\tile)$
almost everywhere with respect to both measures. Hence
\[
\Phi_*\mu(T)=\mu(\tile)=\int_I\sum_{g\in\Z^n}\chi_{T+g}d\lambda=\sum_{g\in\Z^n}\int_{I+g}\chi_T
d\lambda=\int_{\mathbb{R}^n}\chi_T d\lambda=\lambda(T),
\]
where $I$ is the unit cube in $\mathbb{R}^n$. Since $\mu(\tile)$ is positive, $c=1$.
\end{proof}

%To do that, recall that $\mu(\tile)=\mgr(\nucl)$ by the proposition~\ref{theor_tile_measure}. Let
%$K=\Phi(\tile)$. Then $\nu(K)=\mgr(\nucl)$. If $\chi_{K+g}$ is the characteristic function of $K+g$, then
%$\sum_g\chi_{K+g}(t)$ is the number of such $g\in\mathbb{Z}^n$ that $t$ belongs to $K+g$, or, which is the same,
%that $\Phi^{-1}t$ belongs to $\tile\cdot g$. By the proposition~\ref{points_above}
%$\sum_g\chi_{K+g}=\mgr(\nucl)$ almost everywhere with respect to measure $\nu$. It follows that the equality
%holds also a.e. with respect to $\lambda$. Let $I$ be the unit cube in $\mathbb{R}^n$. Then
%$$\mgr(\nucl)=\int_I\sum_g\chi_{K+g}d\lambda=\sum_g\int_{I+g}\chi_K
%d\lambda=\int\chi_K d\lambda=\lambda(K)$$. Thus
%$\nu(K)=\mgr(\nucl)=\lambda(K)$. Since $\tile$ contains interior
%points, its Lebesgue measure is non-zero. Thus $c=1$.


\begin{cor}
The Lebesgue measure of the self-affine tile $T$ is equal to the measure number $\mgr(\nucl)$ of
the nucleus $\nucl$ of the associated self-similar action $(\mathbb{Z}^n,X^{*})$.
\end{cor}

The nucleus of a contracting self-similar action can be found algorithmically using the program
package \cite{AutomGrp}, and the number $\mgr(\nucl)$ can be found using the remarks after
Proposition~\ref{prop_eigenvect_transition_matr}. The measures of sets $T\cap (T+a)$ for
$a\in\mathbb{Z}^n$ can be found by Remark~\ref{rem_meas_inters_tiles}.

%Similar arguments together with Corollary~\ref{cor_meas_gr_rational} can be used to get the result
%of \cite{he_hui:self_affine} that self-affine regions have rational Lebesgue measure.


%\begin{example}
%\red{Some example of computing the self-similar action and the measure of the self-affine sets. Use
%nucleus, the picture from Example in Section~\ref{section_Sofic}}
%\end{example}




The methods developed in \cite{gab_yu:natur_til} to find the Lebesgue measure of integral
self-affine tiles are related to the discussion above. Take the complete automaton
\cite[page~11]{self_sim_groups} of the self-similar action $(\Z^n,X^{*})$ (it actually coincides
with the graph $B(\Z^n)$ from \cite{canon_numsys}), revert the direction of every edge, identify
edges with the same starting and end vertices labeled by $(d_1,d_2)$ with the same difference
$r=d_2-d_1$, and put the new label $r$ on this edge. We get the graph $\mathcal{G}(\Z^n)$
constructed in \cite{gab_yu:natur_til} and \cite{fract_numsys}. The set $W$ constructed using
$\mathcal{G}(\Z^n)$ \cite[page~195]{gab_yu:natur_til} is precisely the nucleus $\nucl$. Hence the
theory of self-similar groups provides a nice explanation to the ideas in
\cite[Section~3]{gab_yu:natur_til} and the methods developed in Sections~2,4 can be seen as its
non-abelian generalization.

It is shown in \cite{lag_wang:int_self_affine_II} that integral self-affine tile $T$ gives a
lattice tiling  of $\mathbb{R}^n$ with some lattice $L\subset\Z^n$. An interesting open question is
whether this holds for any (self-replicating) contracting self-similar action $(G,X^{*})$ (or at
least for self-similar actions of torsion-free nilpotent groups), i.e. the tile $\tile$ gives a
tiling of $\limGs$ with some subgroup $H< G$.



\vspace{0.3cm}\noindent\textbf{Invariant measures on Julia sets of rational functions.} Let
$f(z)\in\mathbb{C}(z)$ be a rational map of the Riemann sphere $\Rs$ of degree $d\geq 2$. It is
proved in \cite{invmeas_rat} that there exists a natural $f$-invariant probability measure $\mu_f$
with the support on the Julia set $J(f)$. In \cite{rational_bernoulli}, D.~Heicklen and C.~Hoffman
showed that the dynamical system $(J(f),f,\mu_f)$ is conjugated to the one-sided Bernoulli
$d$-shift, proving the conjecture in \cite{invmeas_rat,lub:entr_ration}. In the case of
sub-hyperbolic rational functions this result also follows from
Theorem~\ref{theor_lims_conju_Bernoulli} using the theory of iterated monodromy groups developed in
\cite{self_sim_groups}.

Let $f$ be a sub-hyperbolic rational function and let $\M=\Rs\setminus\overline{P}_f$, where $P_f$
is the post-critical set of $f$. Then $f$ defines a $d$-fold partial self-covering $f:
f^{-1}(\M)\rightarrow\M$. Take a base point $t\in\M$ and let $T_t$ be the tree of preimages
$f^{-n}(t)$, $n\geq 0$, where every vertex $z\in f^{-n}(t)$ is connected by an edge to $f(z)\in
f^{-n+1}(t)$. The fundamental group $\pi_1(\M,t)$ acts by automorphisms on $T_t$ through the
monodromy action on every level $f^{-n}(t)$. The quotient of $\pi_1(\M,t)$ by the kernel of its
action on $T_t$ is the \textit{iterated monodromy group} $IMG(f)$ of the map $f$. The group
$IMG(f)$ is contracting self-similar group and the limit dynamical system $(\lims[IMG(f)],\si)$ is
conjugated to $(J(f),f)$ by Theorem~6.4.4 in \cite{self_sim_groups}. Under this conjugation the
measure $\muls$ is send to the measure $\mu_f$, and we can apply
Theorem~\ref{theor_lims_conju_Bernoulli} to the dynamical system $(J(f),f,\mu_f)$.




%Brolin, Hans,  Invariant  sets under  iteration  of  rational functions, Ark. Mat. 6 (1965),
%103-144.

\bibliographystyle{plain}
%\bibliography{bibl}
\begin{thebibliography}{10}

\bibitem{bandt:self_sim_meas}
C.~Bandt, Self-similar measures, {Fiedler, Bernold (ed.), Ergodic theory,
  analysis, and efficient simulation of dynamical systems. Berlin: Springer.
  31-46 (2001).}

\bibitem{fractal_gr_sets}
L.~Bartholdi, R.~Grigorchuk, V.~Nekrashevych, From fractal groups to fractal
  sets, in: Fractals in Graz 2001, Trends Math., Birkh\"auser, Basel, 2003, pp.
  25--118.

\bibitem{branch}
L.~Bartholdi, R.~I. Grigorchuk, Z.~{\v{S}}uni{\'k}, Branch groups, in: Handbook
  of algebra, Vol. 3, North-Holland, Amsterdam, 2003, pp. 989--1112.

\bibitem{sidki:virt_nilp}
A.~Berlatto, S.~Sidki, Virtual endomorphisms of nilpotent groups, Groups
  Geom. Dyn. 1~(1) (2007) 21--46.

\bibitem{deng_he:int_self_affine}
G.-T. Deng, X.-G. He, Integral self-affine sets with positive Lebesgue
  measures, Arch. Math. 90~(2) (2008) 150--157.

\bibitem{edgar:int_prob_fr_mes}
G.~A. Edgar, Integral, probability, and fractal measures, {New York, NY:
  Springer. x, 286 p.}, 1998.

\bibitem{int_two_tiles}
M.~Elekes, T.~Keleti, A.~Ma'the', Self-similar and self-affine sets: measure of the intersection of
two copies. Ergodic Theory and Dynamical Systems,  Published online by Cambridge University Press
29 Jun 2009 doi:10.1017/S0143385709000121


\bibitem{invmeas_rat}
A.~Freire, A.~Lopes, R.~Mane, An invariant measure for rational maps, Bol.
  Soc. Bras. Mat. 14~(1) (1983) 45--62.

\bibitem{gab_yu:natur_til}
J.-P. Gabardo, X.~Yu, Natural tiling, lattice tiling and Lebesgue measure of
  integral self-affine tiles, J. Lond. Math. Soc., II. Ser. 74~(1) (2006)
  184--204.

\bibitem{gri:MilnorProb}
R.~I. Grigorchuk, On the {M}ilnor problem of group growth, Dokl. Akad. Nauk
  SSSR 271~(1) (1983) 30--33.

\bibitem{mult_an_self_sim}
K.~Grochenig and W.R. Madych, Multiresolution analysis, Haar bases, and self-similar tilings of
  $\mathbb{R}^n$,  IEEE Trans. Inf. Theory 38~(2) (1992) 556-568.


\bibitem{he_hui:self_affine}
X.-G. He, K.-S. Lau, H.~Rao, Self-affine sets and graph-directed systems,
  Constructive Approximation 19~(3) (2003) 373--397.

\bibitem{rational_bernoulli}
D.~Heicklen, C.~Hoffman, {Rational maps are $d$-adic Bernoulli}, Ann. of Math.
  (2) 156~(1) (2002) 103--114.

\bibitem{hoffman:Bernoulli}
C.~Hoffman, D.~Rudolph, Uniform endomorphisms which are isomorphic to a
  {B}ernoulli shift, Ann. of Math. (2) 156~(1) (2002) 79--101.

\bibitem{hutch:fract_self_similarity}
J.~E. Hutchinson, {Fractals and self similarity}, Indiana Univ. Math. J. 30
  (1981) 713--747.

\bibitem{lag_wang:haar}
J.~C. Lagarias, Y.~Wang, {Haar type orthonormal wavelet bases in
  $\mathbb{R}^2$}, J. Fourier Anal. Appl. 2~(1) (1995) 1--14.

\bibitem{lag_wang:int_self_affine_I}
J.~C. Lagarias, Y.~Wang, {Integral self-affine tiles in $\mathbb{R}^n$. I:
  Standard and nonstandard digit sets}, J. Lond. Math. Soc., II. Ser. 54~(1)
  (1996) 161--179.

\bibitem{lag_wang:self_affine}
J.~C. Lagarias, Y.~Wang, {Self-affine tiles in $\mathbb{R}^n$}, Adv. Math.
  121~(1) (1996) 21--49.

\bibitem{lag_wang:int_self_affine_II}
J.~C. Lagarias, Y.~Wang, {Integral self-affine tiles in $\mathbb{R}^n$. II:
  Lattice tilings}, J. Fourier Anal. Appl. 3~(1) (1997) 83--102.

\bibitem{ifs_overl_self_sim_mes}
K.-S. Lau, S.-M. Ngai, H.~Rao, {Iterated function systems with overlaps and
  self-similar measures}, J. Lond. Math. Soc., II. Ser. 63~(1) (2001) 99--116.

\bibitem{IntroSymbDyn}
D.~A. Lind, B.~Marcus, An Introduction to Symbolic Dynamics and Coding,
  Cambridge University Press, New York, NY, USA, 1995.

\bibitem{lub:entr_ration}
M.~Lyubich, {Entropy properties of rational endomorphisms of the Riemann
  sphere}, Ergodic Theory Dyn. Syst. 3 (1983) 351--385.

\bibitem{fract_numsys}
W.~Muller, J.~M. Thuswaldner, R.~F. Tichy, {Fractal properties of number
  systems}, Period. Math. Hung. 42~(1-2) (2001) 51--68.


\bibitem{AutomGrp}
Y.~Muntyan and D.~Savchuk.
\newblock {\em {AutomGrp -- GAP package for computations in self-similar groups
  and semigroups, Version 1.1.4.1}}, 2008.


\bibitem{self_sim_groups}
V.~Nekrashevych, Self-similar groups, Vol. 117 of Mathematical Surveys and
  Monographs, American Mathematical Society, Providence, RI, 2005.

\bibitem{handbook:rothblum}
U.~G. Rothblum, Nonnegative matrices and stochastic matrices, in: L.~Hogben
  (Ed.), Handbook of Linear Algebra (Discrete Mathematics and Its
  Applications), {Chapman \& Hall/CRC}, 2006.

\bibitem{canon_numsys}
K.~Scheicher, J.~M. Thuswaldner, {Canonical number systems, counting automata
  and fractals}, Math. Proc. Camb. Philos. Soc. 133~(1) (2002) 163--182.

\bibitem{dyn_sysII}
Y.~E. Sinai, R.~E. Gamkrelidze, {Dynamical systems II. Ergodic theory with
  applications to dynamical systems and statistical mechanics. Transl. from the
  Russian.}, {Encyclopaedia of Mathematical Sciences. 2. Berlin etc.:
  Springer-Verlag. viii, 281 p. }, 1989.

\bibitem{str:self_harm}
R.~S. Strichartz, {Self-similarity in harmonic analysis}, J. Fourier Anal.
  Appl. 1~(1) (1994) 1--37.

\bibitem{urb:meas_dim_confdyn}
M.~Urbanski, {Measures and dimensions in conformal dynamics}, Bull. Am. Math.
  Soc., New Ser. 40~(3) (2003) 281--321.

\bibitem{vince:digit}
A.~Vince, {Digit tiling of Euclidean space}, {Baake, Michael (ed.) et al.,
  Directions in mathematical quasicrystals. Providence, RI: AMS, American
  Mathematical Society. CRM Monogr. Ser. 13, 329-370 (2000).} (2000).


\end{thebibliography}




\end{document}
